\section{Lattice finiteness over the Laurent ring}
\label{pa-chap:Algebra}

This chapter isolates a general finiteness lemma about pairs of lattices in a module over
the ring of Laurent polynomials. Write \(k[T, T^{-1}]\) for the ring of Laurent polynomials
over a commutative ring \(k\) (Chapter~\ref{pa-chap:Curves}): the free \(k\)-module on the
monomials \(T^n\), \(n \in \Z\), with multiplication \(T^m \cdot T^n = T^{m+n}\), so that
\(T^0 = 1\) and \(T^{-1}\) is a two-sided inverse of \(T\). For a subset \(s\) of a
\(k[T,T^{-1}]\)-module \(N\), write \(\langle s \rangle_{k[T,T^{-1}]}\) for the
\(k[T,T^{-1}]\)-submodule it generates. Call a \(k\)-module \emph{finite} if it is finitely
generated over \(k\), and for \(k\)-submodules \(Q_0, Q_1\) of a \(k\)-module, write
\(Q_0 + Q_1\) for the smallest \(k\)-submodule containing both.

The result abstracts the classical computation, for the projective line covered by its two
standard affine charts, of the cohomology of a coherent sheaf as the cokernel of the
difference-of-restrictions map out of the two chart section rings into the ring of sections
on the overlap: two \(k\)-submodules \(Q_0, Q_1\) of a \(k[T,T^{-1}]\)-module \(N\), one
stable under multiplication by \(T\) and containing enough of \(N\), the other stable under
multiplication by \(T^{-1}\) and eventually absorbing every element of \(N\), leave a quotient
\(N / (Q_0 + Q_1)\) that is a finite \(k\)-module — because only a finite window of monomials
can fail to be absorbed by one lattice or the other.

\section{Stability under powers of the generator}

Throughout this section \(N\) is an abelian group carrying both a \(k\)-module structure and a
\(k[T,T^{-1}]\)-module structure (the two need not be related).

\begin{lemma}[Stability under nonnegative powers of \(T\)]
  \label{pa-lem:T_pow_stable}
  \lean{LaurentPolynomial.T_natCast_smul_mem}
  \leanok
  \dcref{custom}

  Let \(Q \subseteq N\) be a \(k\)-submodule such that \(T \cdot x \in Q\) for every
  \(x \in Q\). Then for every \(x \in Q\) and every integer \(m \ge 0\), \(T^m \cdot x \in Q\).
\end{lemma}
\begin{proof}
  \leanok
  Induction on \(m\). For \(m = 0\), \(T^0 \cdot x = x \in Q\). For the step, suppose
  \(T^m \cdot x \in Q\); then
  \[
    T^{m+1} \cdot x = T \cdot (T^m \cdot x) \in Q
  \]
  by the hypothesis on \(Q\) applied to the element \(T^m \cdot x \in Q\).
\end{proof}

\begin{lemma}[Stability under nonnegative powers of \(T^{-1}\)]
  \label{pa-lem:T_inv_pow_stable}
  \lean{LaurentPolynomial.T_neg_natCast_smul_mem}
  \leanok
  \dcref{custom}

  Let \(Q \subseteq N\) be a \(k\)-submodule such that \(T^{-1} \cdot x \in Q\) for every
  \(x \in Q\). Then for every \(x \in Q\) and every integer \(m \ge 0\),
  \(T^{-m} \cdot x \in Q\).
\end{lemma}
\begin{proof}
  \leanok
  Induction on \(m\), exactly as in the proof of Lemma~\ref{pa-lem:T_pow_stable}, with \(T^{-1}\)
  in place of \(T\): for \(m = 0\), \(T^0 \cdot x = x \in Q\); for the step, if
  \(T^{-m} \cdot x \in Q\) then
  \[
    T^{-(m+1)} \cdot x = T^{-1} \cdot (T^{-m} \cdot x) \in Q
  \]
  by the hypothesis on \(Q\).
\end{proof}

\section{The two-lattice lemma}

\begin{theorem}[Two-lattice finiteness, spanning-set form]
  \label{pa-thm:two_lattice_span}
  \lean{LaurentPolynomial.moduleFinite_quotient_sup_of_span_eq_top}
  \leanok
  \source{hartshorne-algebraic-geometry:page-0242, hartshorne-algebraic-geometry:page-0243}

  \dcref{custom}
  Let \(k\) be a commutative ring and \(N\) a module over \(k[T,T^{-1}]\) whose
  \(k\)-module structure is the restriction of scalars of the \(k[T,T^{-1}]\)-action along the
  canonical ring map \(k \to k[T,T^{-1}]\). Let \(Q_0, Q_1 \subseteq N\) be \(k\)-submodules
  and \(s \subset N\) a finite subset such that:
  \begin{enumerate}
    \item \(s \subseteq Q_0\);
    \item \(s\) spans \(N\) over \(k[T,T^{-1}]\): \(\langle s \rangle_{k[T,T^{-1}]} = N\);
    \item \(Q_0\) is stable under multiplication by \(T\): \(T \cdot x \in Q_0\) for all
      \(x \in Q_0\);
    \item \(Q_1\) is stable under multiplication by \(T^{-1}\): \(T^{-1} \cdot x \in Q_1\) for
      all \(x \in Q_1\);
    \item every \(n \in N\) is eventually absorbed by \(Q_1\): there is an integer \(m \ge 0\)
      with \(T^{-m} \cdot n \in Q_1\).
  \end{enumerate}
  Then \(N / (Q_0 + Q_1)\) is a finite \(k\)-module: it is spanned by the classes of the
  monomials \(T^j \cdot x\), for \(x \in s\) and \(-M < j < 0\), where \(M\) is any uniform
  bound with \(T^{-M} \cdot x \in Q_1\) for every \(x \in s\).
\end{theorem}
\begin{proof}
  \leanok
  \uses{pa-lem:T_pow_stable, pa-lem:T_inv_pow_stable}
  Since \(s\) is finite, hypothesis (5) applied to each \(x \in s\) gives an integer
  \(m_x \ge 0\) with \(T^{-m_x} \cdot x \in Q_1\); let \(M = \max_{x \in s} m_x\) (or \(M = 0\)
  if \(s\) is empty). For each \(x \in s\), since \(M - m_x \ge 0\) and \(Q_1\) is stable
  under \(T^{-1}\), Lemma~\ref{pa-lem:T_inv_pow_stable} gives
  \[
    T^{-M} \cdot x = T^{-(M - m_x)} \cdot \bigl(T^{-m_x} \cdot x\bigr) \in Q_1 ,
  \]
  so \(M\) is a uniform bound for all of \(s\).

  Let \(F = \{\, T^j \cdot x : x \in s,\ -M < j < 0 \,\}\), a finite subset of \(N\) (finite
  since \(s\) is finite and \(j\) ranges over the \(M - 1\) integers strictly between \(-M\)
  and \(0\)).

  \emph{Claim.} For every \(x \in s\) and every \(j \in \Z\),
  \(T^j \cdot x \in Q_0 + Q_1 + \langle F \rangle_k\), where \(\langle F \rangle_k\) is the
  \(k\)-submodule generated by \(F\).

  Indeed, if \(j \ge 0\), then since \(x \in s \subseteq Q_0\) and \(Q_0\) is stable under
  \(T\), Lemma~\ref{pa-lem:T_pow_stable} gives \(T^j \cdot x \in Q_0\). If \(j \le -M\), write
  \(j = -d - M\) with \(d = -j - M \ge 0\); then
  \[
    T^j \cdot x = T^{-d} \cdot \bigl(T^{-M} \cdot x\bigr) ,
  \]
  and since \(T^{-M} \cdot x \in Q_1\) and \(Q_1\) is stable under \(T^{-1}\),
  Lemma~\ref{pa-lem:T_inv_pow_stable} gives \(T^j \cdot x \in Q_1\). Finally, if \(-M < j < 0\),
  then \(T^j \cdot x \in F \subseteq \langle F \rangle_k\) by definition of \(F\). This proves
  the claim.

  Every element \(f \in k[T,T^{-1}]\) is, by definition of the Laurent polynomial ring, a
  \(k\)-linear combination of finitely many monomials \(T^j\), \(j \in \Z\). Since
  \(Q_0 + Q_1 + \langle F \rangle_k\) is a \(k\)-submodule of \(N\), the claim and this
  decomposition of \(f\) give, for every \(x \in s\),
  \[
    f \cdot x \in Q_0 + Q_1 + \langle F \rangle_k .
  \]
  Now let \(n \in N\) be arbitrary. By hypothesis (2), \(n\) is a (finite) \(k[T,T^{-1}]\)-
  linear combination \(n = \sum_{x \in s} f_x \cdot x\) of elements of \(s\); each summand lies
  in \(Q_0 + Q_1 + \langle F \rangle_k\) by the previous paragraph, and this submodule is
  closed under addition, so \(n \in Q_0 + Q_1 + \langle F \rangle_k\). Hence
  \[
    Q_0 + Q_1 + \langle F \rangle_k = N .
  \]
  Consequently the composite \(\langle F \rangle_k \hookrightarrow N \twoheadrightarrow
  N/(Q_0+Q_1)\) is surjective, so the (finite) image of \(F\) in \(N/(Q_0+Q_1)\) spans it as a
  \(k\)-module. Thus \(N/(Q_0+Q_1)\) is a finite \(k\)-module.
\end{proof}

\section{Localization forms}

The spanning-set form of Theorem~\ref{pa-thm:two_lattice_span} is inconvenient when only a
localization-type hypothesis on \(Q_0\) is available, symmetric to hypothesis (5) on \(Q_1\):
every element of \(N\) is eventually absorbed by \(Q_0\) after multiplying by a sufficiently
large power of \(T\). The next two results derive the finiteness conclusion from this
symmetric pair of hypotheses, at the cost of assuming \(N\) itself is a finite
\(k[T,T^{-1}]\)-module; this is exactly the shape of hypothesis produced by localizing a
finite module away from an element.

\begin{corollary}[Two-lattice finiteness, localization form]
  \label{pa-cor:two_lattice_loc}
  \lean{LaurentPolynomial.moduleFinite_quotient_sup_of_exists_smul_mem}
  \leanok
  \dcref{custom}

  Let \(k\) be a commutative ring and \(N\) a finite module over \(k[T,T^{-1}]\) whose
  \(k\)-module structure is the restriction of scalars of the \(k[T,T^{-1}]\)-action along the
  canonical ring map \(k \to k[T,T^{-1}]\). Let \(Q_0, Q_1 \subseteq N\) be \(k\)-submodules
  such that:
  \begin{enumerate}
    \item \(Q_0\) is stable under multiplication by \(T\);
    \item \(Q_1\) is stable under multiplication by \(T^{-1}\);
    \item every \(n \in N\) is eventually absorbed by \(Q_0\): there is an integer \(m \ge 0\)
      with \(T^m \cdot n \in Q_0\);
    \item every \(n \in N\) is eventually absorbed by \(Q_1\): there is an integer \(m \ge 0\)
      with \(T^{-m} \cdot n \in Q_1\).
  \end{enumerate}
  Then \(N / (Q_0 + Q_1)\) is a finite \(k\)-module.
\end{corollary}
\begin{proof}
  \leanok
  \uses{pa-thm:two_lattice_span}
  Since \(N\) is a finite \(k[T,T^{-1}]\)-module, choose a finite subset \(t \subset N\) with
  \(\langle t \rangle_{k[T,T^{-1}]} = N\). By hypothesis (3), for each \(y \in t\) choose an
  integer \(\mu(y) \ge 0\) with \(T^{\mu(y)} \cdot y \in Q_0\), and set
  \[
    s = \{\, T^{\mu(y)} \cdot y : y \in t \,\} \subseteq Q_0 ,
  \]
  a finite subset of \(N\).

  The set \(s\) again spans \(N\) over \(k[T,T^{-1}]\): each element of \(s\) is a
  \(k[T,T^{-1}]\)-multiple of an element of \(t\), so
  \(\langle s \rangle_{k[T,T^{-1}]} \subseteq \langle t \rangle_{k[T,T^{-1}]} = N\); and since
  \(T\) is invertible in \(k[T,T^{-1}]\) with inverse \(T^{-1}\), every \(y \in t\) is
  recovered from \(s\) by
  \[
    y = T^{-\mu(y)} \cdot \bigl(T^{\mu(y)} \cdot y\bigr) ,
  \]
  a \(k[T,T^{-1}]\)-multiple of the element \(T^{\mu(y)} \cdot y \in s\); so
  \(t \subseteq \langle s \rangle_{k[T,T^{-1}]}\), giving
  \(\langle s \rangle_{k[T,T^{-1}]} \supseteq \langle t \rangle_{k[T,T^{-1}]} = N\) as well.
  Hence \(\langle s \rangle_{k[T,T^{-1}]} = N\).

  Now \(s \subseteq Q_0\), \(s\) spans \(N\) over \(k[T,T^{-1}]\), \(Q_0\) is \(T\)-stable,
  \(Q_1\) is \(T^{-1}\)-stable, and hypothesis (4) is exactly hypothesis (5) of
  Theorem~\ref{pa-thm:two_lattice_span}. Applying that theorem to \(s\), \(Q_0\), \(Q_1\) gives
  that \(N/(Q_0+Q_1)\) is a finite \(k\)-module.
\end{proof}

\begin{corollary}[Two-lattice finiteness, localization form via powers of the generator]
  \label{pa-cor:two_lattice_loc_pow}
  \lean{LaurentPolynomial.moduleFinite_quotient_sup_of_exists_pow_smul_mem}
  \leanok
  \dcref{custom}

  Let \(k\), \(N\), \(Q_0\), \(Q_1\) be as in Corollary~\ref{pa-cor:two_lattice_loc}, satisfying
  hypotheses (1) and (2) there, and suppose in place of (3) and (4) that:
  \begin{enumerate}
    \item[(3a)] every \(n \in N\) satisfies: there is an integer \(m \ge 0\) with
      \(T^m \cdot n \in Q_0\), where \(T^m\) denotes the ordinary \(m\)-th power of the ring
      element \(T \in k[T,T^{-1}]\);
    \item[(4a)] every \(n \in N\) satisfies: there is an integer \(m \ge 0\) with
      \(T^{-m} \cdot n \in Q_1\), where \(T^{-m}\) denotes the ordinary \(m\)-th power of
      \(T^{-1} \in k[T,T^{-1}]\).
  \end{enumerate}
  Then \(N / (Q_0 + Q_1)\) is a finite \(k\)-module.
\end{corollary}
\begin{proof}
  \leanok
  \uses{pa-cor:two_lattice_loc}
  For every integer \(m \ge 0\), the \(m\)-th power of the ring element \(T\) is the monomial
  \(T^m\), and the \(m\)-th power of \(T^{-1}\) is the monomial \(T^{-m}\); so hypotheses
  (3a) and (4a) here coincide, element for element, with hypotheses (3) and (4) of
  Corollary~\ref{pa-cor:two_lattice_loc}. The conclusion is immediate from that corollary.
\end{proof}

\section{Tensor products of localizations and the descent decompositions}
\label{pa-sec:tensorAway}

This section collects two pieces of commutative-algebra scaffolding for the descent along a
finite localization cover. Let \(A \to B\) be a homomorphism of commutative rings. Recall that
for \(r \in B\), a \emph{model of the localization of \(B\) away from \(r\)} is a \(B\)-algebra
\(S\) in which the image of \(r\) is invertible and which is universal with this property,
so that \(S \cong B[1/r]\); we treat such a model as an \(A\)-algebra through the composite
\(A \to B \to S\). The first part exhibits the \(A\)-tensor product of two such localization
models — taken over two possibly distinct bases \(B_1, B_2\), each a commutative
\(A\)-algebra — as itself a localization; the second decomposes the tensor square and cube of a
finite product of localizations into their pairwise and triple components.

\subsection{The tensor product of two localizations}

\begin{definition}[The tensor-localization algebra structure]
  \label{pa-def:tensorAwayAlgebra}
  \lean{IsLocalization.Away.tensorMap, IsLocalization.Away.tensorAwayAlgebra}
  \leanok
  \dcref{custom}

  Let \(A \to B_1\) and \(A \to B_2\) be homomorphisms of commutative rings, let \(r \in B_1\)
  and \(s \in B_2\), and let \(S_i, S_j\) be models of the localization of \(B_1\) away from
  \(r\), respectively of \(B_2\) away from \(s\), regarded as \(A\)-algebras through
  \(A \to B_1 \to S_i\) and \(A \to B_2 \to S_j\). Tensoring the two structure maps
  \(B_1 \to S_i\) and \(B_2 \to S_j\) over \(A\) gives an \(A\)-algebra homomorphism
  \[
    B_1 \otimes_A B_2 \longrightarrow S_i \otimes_A S_j , \qquad
    b_1 \otimes b_2 \mapsto \bar b_1 \otimes \bar b_2 ,
  \]
  whose underlying ring map endows \(S_i \otimes_A S_j\) with the structure of a
  \((B_1 \otimes_A B_2)\)-algebra.
\end{definition}

\begin{theorem}[The tensor product of two localizations is a localization of the tensor square]
  \label{pa-thm:isLocalization_away_tensor}
  \lean{IsLocalization.Away.isLocalization_away_tensor}
  \leanok
  \dcref{custom}

  \uses{pa-def:tensorAwayAlgebra}
  With the structure of Definition~\ref{pa-def:tensorAwayAlgebra}, \(S_i \otimes_A S_j\) is the
  localization of \(B_1 \otimes_A B_2\) away from the element
  \((r \otimes 1)(1 \otimes s) = r \otimes s\).
\end{theorem}
\begin{proof}
  \leanok
  Factor the structure map \(B_1 \otimes_A B_2 \to S_i \otimes_A S_j\) through the intermediate
  ring \(S_i \otimes_A B_2\). The first stage \(B_1 \otimes_A B_2 \to S_i \otimes_A B_2\)
  localizes only the left factor: it is the base change of the localization \(B_1 \to S_i\) at
  \(r\) along the ring map \(B_1 \to B_1 \otimes_A B_2\), hence exhibits \(S_i \otimes_A B_2\) as
  the localization of \(B_1 \otimes_A B_2\) away from \(r \otimes 1\). The second stage
  \(S_i \otimes_A B_2 \to S_i \otimes_A S_j\) localizes the right factor: it is the base change of
  the localization \(B_2 \to S_j\) at \(s\) along \(B_2 \to S_i \otimes_A B_2\), exhibiting
  \(S_i \otimes_A S_j\) as the localization of \(S_i \otimes_A B_2\) away from \(1 \otimes s\).
  Two successive away-localizations, first at \(r \otimes 1\) then at \(1 \otimes s\), compose
  to the away-localization at the product \((r \otimes 1)(1 \otimes s)\); and by the
  multiplication rule in \(B_1 \otimes_A B_2\) this product equals \(r \otimes s\).
\end{proof}

\begin{definition}[Comparison with an arbitrary model]
  \label{pa-def:tensorAwayEquiv}
  \lean{IsLocalization.Away.tensorAwayEquiv}
  \leanok
  \dcref{custom}

  \uses{pa-thm:isLocalization_away_tensor}
  Let \(T\) be any model of the localization of \(B_1 \otimes_A B_2\) away from
  \((r \otimes 1)(1 \otimes s)\). The canonical comparison of two localizations of
  \(B_1 \otimes_A B_2\) at the same element is a \((B_1 \otimes_A B_2)\)-algebra isomorphism
  \[
    S_i \otimes_A S_j \;\xrightarrow{\ \sim\ }\; T .
  \]
\end{definition}

\begin{lemma}[Values of the comparison isomorphism]
  \label{pa-lem:tensorAwayEquiv_tmul}
  \lean{IsLocalization.Away.tensorAwayEquiv_tensorMap, IsLocalization.Away.tensorAwayEquiv_tmul}
  \leanok
  \dcref{custom}

  \uses{pa-def:tensorAwayEquiv}
  The comparison isomorphism sends the image of \(w \in B_1 \otimes_A B_2\) under the structure
  map to the image of \(w\) in \(T\); in particular, on a pure tensor of structure-map images it
  sends \(\bar b_1 \otimes \bar b_2\) to the image of \(b_1 \otimes b_2\) in \(T\).
\end{lemma}
\begin{proof}
  \leanok
  Being a \((B \otimes_A B)\)-algebra homomorphism, the comparison isomorphism commutes with the
  two structure maps, which is the first statement. The second follows because the pure tensor
  \(\bar b_1 \otimes \bar b_2\) is the structure-map image of \(b_1 \otimes b_2\).
\end{proof}

\subsection{Decompositions of the tensor square and cube of a finite product}

Throughout this subsection \(A\) is a commutative ring, \(\iota\) a finite index set, and
\((S_i)_{i \in \iota}\) a finite family of \(A\)-algebras with product \(P := \prod_{i} S_i\).

\begin{definition}[The pairwise and triple decompositions]
  \label{pa-def:piDoubleEquivA}
  \lean{Algebra.TensorProduct.piDoubleEquivA, Algebra.TensorProduct.piDoubleEquivA_tmul}
  \leanok
  \dcref{custom}

  Distributing the tensor product over the finite product in each factor gives the canonical
  \(A\)-algebra identification
  \[
    P \otimes_A P \;\cong\; \prod_{(i,j) \in \iota \times \iota} S_i \otimes_A S_j ,
  \]
  which sends a pure tensor \(s \otimes t\) to the family \((i,j) \mapsto s_i \otimes t_j\).
\end{definition}

\begin{definition}[The triple decomposition]
  \label{pa-def:piTripleEquivA}
  \lean{Algebra.TensorProduct.piTripleEquivA, Algebra.TensorProduct.piTripleEquivA_tmul}
  \leanok
  \dcref{custom}

  \uses{pa-def:piDoubleEquivA}
  Applying the pairwise decomposition to the inner factor and distributing once more gives the
  canonical \(A\)-algebra identification
  \[
    P \otimes_A (P \otimes_A P) \;\cong\;
      \prod_{(i,j,k) \in \iota^3} S_i \otimes_A (S_j \otimes_A S_k) ,
  \]
  sending \(s \otimes v\) to the family \((i,j,k) \mapsto s_i \otimes (v_{(j,k)})\), where
  \(v_{(j,k)}\) is the \((j,k)\)-component of \(v\) under the pairwise decomposition.
\end{definition}

\begin{definition}[The index-wise simplicial maps]
  \label{pa-def:indexwiseMaps}
  \lean{Algebra.TensorProduct.inclLeftA, Algebra.TensorProduct.inclRightA,
    Algebra.TensorProduct.faceA₁₂, Algebra.TensorProduct.faceA₁₃, Algebra.TensorProduct.faceA₂₃,
    Algebra.TensorProduct.diagA}
  \leanok
  \dcref{custom}

  On the pairwise tensor components one has the index-wise simplicial maps: for indices
  \(i, j, k\), the two \emph{inclusions} \(S_i \to S_i \otimes_A S_j\), \(x \mapsto x \otimes 1\),
  and \(S_j \to S_i \otimes_A S_j\), \(y \mapsto 1 \otimes y\); the three \emph{faces} into the
  triple tensor \(S_i \otimes_A (S_j \otimes_A S_k)\), hitting the positions \(\{1,2\}\),
  \(\{1,3\}\), \(\{2,3\}\) out of \((S_i \otimes_A S_j)\), \((S_i \otimes_A S_k)\),
  \((S_j \otimes_A S_k)\) respectively; and the \emph{diagonal}
  \(S_i \otimes_A S_i \to S_i\), \(x \otimes y \mapsto xy\).
\end{definition}

\begin{lemma}[The descent inclusions become the index-wise inclusions]
  \label{pa-lem:descentIncl_transport}
  \lean{Algebra.TensorProduct.piDoubleEquivA_descentIncl₁,
    Algebra.TensorProduct.piDoubleEquivA_descentIncl₂}
  \leanok
  \dcref{custom}

  \uses{pa-def:piDoubleEquivA, pa-def:indexwiseMaps}
  Under the pairwise decomposition, the two inclusions \(P \to P \otimes_A P\) become the
  index-wise inclusions: \(s \otimes 1\) corresponds to the family
  \((i,j) \mapsto s_i \otimes 1\), and \(1 \otimes s\) to \((i,j) \mapsto 1 \otimes s_j\).
\end{lemma}
\begin{proof}
  \leanok
  Evaluate at a component \((i,j)\): by the pure-tensor formula of
  Definition~\ref{pa-def:piDoubleEquivA}, the decomposition sends \(s \otimes 1\) to
  \(s_i \otimes 1_j = s_i \otimes 1\), which is the first index-wise inclusion of \(s_i\), and
  \(1 \otimes s\) to \(1 \otimes s_j\), the second index-wise inclusion of \(s_j\).
\end{proof}

\begin{lemma}[The descent faces become the index-wise faces]
  \label{pa-lem:descentFace_transport}
  \lean{Algebra.TensorProduct.piTripleEquivA_descentFace₂₃,
    Algebra.TensorProduct.piTripleEquivA_descentFace₁₂,
    Algebra.TensorProduct.piTripleEquivA_descentFace₁₃}
  \leanok
  \dcref{custom}

  \uses{pa-def:piTripleEquivA, pa-def:piDoubleEquivA, pa-def:indexwiseMaps}
  Under the pairwise and triple decompositions, the three faces \(P \otimes_A P \to
  P \otimes_A (P \otimes_A P)\) become the index-wise faces: the \(\{2,3\}\)-face carries a
  family \((w_{(i,j)})\) to \((i,j,k) \mapsto \mathrm{face}_{23}(w_{(j,k)})\), and likewise the
  \(\{1,2\}\)- and \(\{1,3\}\)-faces carry it to
  \((i,j,k) \mapsto \mathrm{face}_{12}(w_{(i,j)})\) and
  \((i,j,k) \mapsto \mathrm{face}_{13}(w_{(i,k)})\).
\end{lemma}
\begin{proof}
  \leanok
  Each identity is checked on generators using the pure-tensor formulas of
  Definitions~\ref{pa-def:piDoubleEquivA} and~\ref{pa-def:piTripleEquivA}. For the \(\{2,3\}\)-face,
  the source element \(1 \otimes v\) is carried to
  \((i,j,k) \mapsto 1 \otimes (v_{(j,k)}) = \mathrm{face}_{23}(w_{(j,k)})\) directly. The
  \(\{1,2\}\)- and \(\{1,3\}\)-faces mix the two product factors; since both sides are
  \(A\)-algebra homomorphisms it suffices to compare them on pure tensors, where the two
  decompositions evaluate to the stated index-wise faces.
\end{proof}

\begin{lemma}[Multiplication becomes the index-wise diagonal]
  \label{pa-lem:lmul_transport}
  \lean{Algebra.TensorProduct.lmul'_piDoubleEquivA_symm}
  \leanok
  \dcref{custom}

  \uses{pa-def:piDoubleEquivA, pa-def:indexwiseMaps}
  Under the pairwise decomposition, the multiplication map \(P \otimes_A P \to P\),
  \(a \otimes b \mapsto ab\), becomes the index-wise diagonal: a family \((w_{(i,j)})\) is
  carried to \(i \mapsto \mathrm{diag}(w_{(i,i)})\), the product of the diagonal components.
\end{lemma}
\begin{proof}
  \leanok
  Both sides are \(A\)-algebra homomorphisms, so it suffices to compare on a pure tensor
  \(s \otimes t\). Multiplication sends it to \(i \mapsto s_i t_i\); the decomposition sends it
  to \((i,j) \mapsto s_i \otimes t_j\), whose diagonal \((i,i)\)-component is \(s_i \otimes t_i\),
  and \(\mathrm{diag}(s_i \otimes t_i) = s_i t_i\). The two agree.
\end{proof}

\begin{theorem}[Faithful flatness of the composite cover tower]
  \label{pa-thm:pi_tower}
  \lean{Module.FaithfullyFlat.pi_tower}
  \leanok
  \dcref{custom}

  Let \(A \to B\) be faithfully flat, let \(f \colon \iota \to B\) be a covering family — the
  \(f_i\) generate the unit ideal of \(B\) — and let each \(S_i\) be a model of the localization
  of \(B\) away from \(f_i\). Then \(P = \prod_i S_i\) is a faithfully flat \(A\)-algebra.
\end{theorem}
\begin{proof}
  \leanok
  The product of the localizations at a covering family is faithfully flat over \(B\): each
  \(B \to S_i\) is flat, being a localization, and because the \(f_i\) generate the unit ideal
  the family jointly detects triviality, so \(B \to \prod_i S_i\) is faithfully flat. Faithful
  flatness is preserved by composition, so composing with the faithfully flat \(A \to B\) makes
  the tower \(A \to B \to \prod_i S_i\) faithfully flat; hence \(P\) is a faithfully flat
  \(A\)-algebra.
\end{proof}

\section{The composite descent unit of a coherent family of units}
\label{pa-sec:piAssembly}

Keep the notation of the previous section: \(A\) is a commutative ring, \(\iota\) a finite
index set, and \((S_i)_{i \in \iota}\) a finite family of \(A\)-algebras with product
\(P := \prod_i S_i\). This section assembles a family of componentwise units, one on each
pairwise tensor \(S_i \otimes_A S_j\), into a single unit of \(P \otimes_A P\), and isolates
the conditions under which that unit is an \(A\)-descent \(1\)-cocycle
(Definition~\ref{pa-def:module_descent_cocycle}) whose collapse
(Definition~\ref{pa-def:module_tensorCollapse}) recovers a prescribed cover cocycle. This is
the pure-algebra core of the composite descent unit assembled by the \'etale-separatedness
campaign.

\begin{definition}[The composite descent unit]
  \label{pa-def:piAssemblyUnit}
  \lean{Algebra.TensorProduct.piAssemblyUnit}
  \leanok
  \dcref{custom}

  \uses{pa-def:piDoubleEquivA}
  Given a family of units \(w_{ij} \in (S_i \otimes_A S_j)^{\times}\), the \emph{composite
  descent unit} \(\mathrm{piAssembly}(w) \in (P \otimes_A P)^{\times}\) is the unit
  corresponding to the family \((w_{ij})\) under the pairwise decomposition
  \(P \otimes_A P \cong \prod_{(i,j)} S_i \otimes_A S_j\) of
  Definition~\ref{pa-def:piDoubleEquivA}.
\end{definition}

\begin{theorem}[The assembled unit of a coherent family is a descent cocycle]
  \label{pa-thm:isDescentCocycle_piAssemblyUnit}
  \lean{Algebra.TensorProduct.isDescentCocycle_piAssemblyUnit}
  \leanok
  \dcref{custom}

  \uses{pa-def:piAssemblyUnit, pa-def:module_descent_cocycle, pa-def:indexwiseMaps,
    pa-def:piTripleEquivA, pa-lem:lmul_transport, pa-lem:descentFace_transport}
  Suppose the family \((w_{ij})\) is:
  \begin{enumerate}
    \item \emph{diagonally normalized}: \(\mathrm{diag}(w_{ii}) = 1\) in \(S_i\) for every
      \(i\), where \(\mathrm{diag} \colon S_i \otimes_A S_i \to S_i\) is the multiplication
      (Definition~\ref{pa-def:indexwiseMaps});
    \item \emph{index-wise cocyclic}: for all \(i, j, k\),
      \(\mathrm{face}_{23}(w_{jk}) \cdot \mathrm{face}_{12}(w_{ij}) =
      \mathrm{face}_{13}(w_{ik})\) in \(S_i \otimes_A (S_j \otimes_A S_k)\).
  \end{enumerate}
  Then the composite descent unit \(\mathrm{piAssembly}(w)\) is an \(A\)-descent
  \(1\)-cocycle relative to \(A \to P\).
\end{theorem}
\begin{proof}
  \leanok
  A descent \(1\)-cocycle is a unit satisfying the normalization \(\mu(u) = 1\) and the
  cocycle identity \(\partial_{23}(u) \cdot \partial_{12}(u) = \partial_{13}(u)\)
  (Definition~\ref{pa-def:module_descent_cocycle}); both conditions are verified after
  transporting through the pairwise and triple decompositions.

  For the normalization: under the pairwise decomposition the multiplication
  \(\mu \colon P \otimes_A P \to P\) becomes the index-wise diagonal
  (Lemma~\ref{pa-lem:lmul_transport}), carrying \(\mathrm{piAssembly}(w)\) to
  \(i \mapsto \mathrm{diag}(w_{ii})\), which is \(1\) by hypothesis~(1); since the
  decomposition is an isomorphism, \(\mu(\mathrm{piAssembly}(w)) = 1\).

  For the cocycle identity: the triple decomposition (Definition~\ref{pa-def:piTripleEquivA})
  is injective, so it suffices to compare the two sides after applying it. Under the two
  decompositions the three faces \(\partial_{23}, \partial_{12}, \partial_{13}\) become the
  index-wise faces \(\mathrm{face}_{23}, \mathrm{face}_{12}, \mathrm{face}_{13}\)
  (Lemma~\ref{pa-lem:descentFace_transport}), so the transported left side is the family
  \((i,j,k) \mapsto \mathrm{face}_{23}(w_{jk}) \cdot \mathrm{face}_{12}(w_{ij})\) and the
  transported right side is \((i,j,k) \mapsto \mathrm{face}_{13}(w_{ik})\); these agree
  componentwise by hypothesis~(2).
\end{proof}

For the remaining results fix, in addition, a commutative \(A\)-algebra \(B\), a finite
family \(f \colon \iota \to B\), models \(S_i\) of the localization of \(B\) away from
\(f_i\), and \emph{overlap models} \(T_{ij}\) of the localization of \(B\) away from
\(f_i f_j\), all regarded as \(A\)-algebras through \(A \to B\). Write
\(\iota_1 \colon S_i \to T_{ij}\) and \(\iota_2 \colon S_j \to T_{ij}\) for the two canonical
\(B\)-algebra structure maps of the overlap model (\(f_i\), resp.\ \(f_j\), being invertible
in \(T_{ij}\)).

\begin{definition}[The index-wise collapse]
  \label{pa-def:componentCollapse}
  \lean{IsLocalization.AwayCover.componentCollapse}
  \leanok
  \dcref{custom}

  \uses{pa-thm:isLocalization_away_tensor}
  The \emph{index-wise collapse} is the \(A\)-algebra homomorphism
  \[
    \kappa_{ij} \colon S_i \otimes_A S_j \longrightarrow T_{ij}, \qquad
    x \otimes y \mapsto \iota_1(x) \cdot \iota_2(y) .
  \]
  It is the \((B \otimes_A B)\)-semilinear comparison of the two-base tensor localization
  \(S_i \otimes_A S_j\) of \(B \otimes_A B\) (Theorem~\ref{pa-thm:isLocalization_away_tensor}
  with \(B_1 = B_2 = B\)) with the one-base overlap localization \(T_{ij}\) of \(B\).
\end{definition}

\begin{lemma}[The collapse computes componentwise]
  \label{pa-lem:piDoubleEquiv_tensorCollapse}
  \lean{IsLocalization.AwayCover.piDoubleEquiv_tensorCollapse}
  \leanok
  \dcref{custom}

  \uses{pa-def:componentCollapse, pa-def:piDoubleEquivA, pa-def:module_tensorCollapse}
  Decompose both \(P \otimes_A P \cong \prod_{ij} S_i \otimes_A S_j\)
  (Definition~\ref{pa-def:piDoubleEquivA}) and the target
  \(P \otimes_B P \cong \prod_{ij} T_{ij}\) of the collapse pairwise. Then the collapse map
  \(\pi \colon P \otimes_A P \to P \otimes_B P\) of
  Definition~\ref{pa-def:module_tensorCollapse} becomes the index-wise collapse: it carries a
  family \((v_{ij})\) to \((i,j) \mapsto \kappa_{ij}(v_{ij})\).
\end{lemma}
\begin{proof}
  \leanok
  Both composites are \(A\)-algebra homomorphisms out of \(P \otimes_A P\), so it suffices to
  compare them on a pure tensor \(s \otimes t\). The pairwise decomposition sends
  \(s \otimes t\) to the family \((i,j) \mapsto s_i \otimes t_j\); applying \(\kappa_{ij}\)
  yields \(\iota_1(s_i) \cdot \iota_2(t_j)\). On the other side, the collapse sends
  \(s \otimes t\) to the same elementary tensor in \(P \otimes_B P\), whose \((i,j)\)-component
  under the \(B\)-side pairwise decomposition is the image of \(s_i \otimes t_j\) under the
  canonical comparison of \(S_i \otimes_B S_j\) with \(T_{ij}\); by uniqueness of \(B\)-algebra
  maps out of a localization of \(B\) this comparison equals \(\kappa_{ij}\) on pure tensors,
  again \(\iota_1(s_i) \cdot \iota_2(t_j)\). The two families coincide.
\end{proof}

\begin{theorem}[The collapse of the assembled unit is the descent unit of the cover cocycle]
  \label{pa-thm:tensorCollapse_piAssemblyUnit}
  \lean{IsLocalization.AwayCover.tensorCollapse_piAssemblyUnit}
  \leanok
  \dcref{custom}

  \uses{pa-def:piAssemblyUnit, pa-def:componentCollapse, pa-lem:piDoubleEquiv_tensorCollapse,
    pa-def:module_tensorCollapse, pa-def:module_descent_cocycle}
  Let \((c_{ij})\), \(c_{ij} \in T_{ij}^{\times}\), be a cover cocycle for \(f\), and suppose
  the componentwise units \(w_{ij}\) collapse onto it: \(\kappa_{ij}(w_{ij}) = c_{ij}\) for
  all \(i, j\). Then the collapse of the composite descent unit along the tower
  \(A \to B \to P\) is the descent unit of \(c\):
  \[
    \pi\bigl(\mathrm{piAssembly}(w)\bigr) = \mathrm{cocycleUnit}(c)
    \quad\text{in } (P \otimes_B P)^{\times},
  \]
  where \(\mathrm{cocycleUnit}(c)\) is the unit of \(P \otimes_B P\) corresponding to the
  family \((c_{ij})\) under the pairwise decomposition
  \(P \otimes_B P \cong \prod_{ij} T_{ij}\).
\end{theorem}
\begin{proof}
  \leanok
  Both sides are units of \(P \otimes_B P\); apply the injective pairwise decomposition
  \(P \otimes_B P \cong \prod_{ij} T_{ij}\). By Lemma~\ref{pa-lem:piDoubleEquiv_tensorCollapse}
  the decomposition of \(\pi(\mathrm{piAssembly}(w))\) is the family
  \((i,j) \mapsto \kappa_{ij}(u_{ij})\), where \(u_{ij}\) is the \((i,j)\)-component of
  \(\mathrm{piAssembly}(w)\); by Definition~\ref{pa-def:piAssemblyUnit} that component is
  \(w_{ij}\), so the family is \((i,j) \mapsto \kappa_{ij}(w_{ij}) = c_{ij}\) by hypothesis.
  The decomposition of \(\mathrm{cocycleUnit}(c)\) is the family \((c_{ij})\) by definition.
  The two families agree, hence so do the two units.
\end{proof}

\begin{theorem}[Diagonal normalization from a collapse onto a normalized family]
  \label{pa-thm:diagA_eq_one_of_componentCollapse}
  \lean{IsLocalization.AwayCover.diagA_eq_one_of_componentCollapse}
  \leanok
  \dcref{custom}

  \uses{pa-def:componentCollapse, pa-def:indexwiseMaps}
  Suppose the componentwise units collapse onto a family \((c_{ij})\),
  \(\kappa_{ij}(w_{ij}) = c_{ij}\), whose diagonal is normalized: the image of \(c_{ii}\)
  under the canonical \(B\)-algebra map \(\delta_i \colon T_{ii} \to S_i\) (well-defined
  because \(f_i\) is already invertible in \(S_i\)) is \(1\), for every \(i\). Then
  \(\mathrm{diag}(w_{ii}) = 1\) in \(S_i\) for every \(i\).
\end{theorem}
\begin{proof}
  \leanok
  Fix \(i\). The composites \(\delta_i \circ \iota_1\) and \(\delta_i \circ \iota_2\) are
  \(B\)-algebra endomorphisms of the localization \(S_i\), hence both equal the identity by
  uniqueness of \(B\)-algebra maps out of \(S_i\). Consequently, for a pure tensor
  \(x \otimes y \in S_i \otimes_A S_i\),
  \[
    \delta_i\bigl(\kappa_{ii}(x \otimes y)\bigr)
      = \delta_i(\iota_1 x) \cdot \delta_i(\iota_2 y) = x \cdot y
      = \mathrm{diag}(x \otimes y) ,
  \]
  and both \(\delta_i \circ \kappa_{ii}\) and \(\mathrm{diag}\) are additive, so
  \(\delta_i \circ \kappa_{ii} = \mathrm{diag}\) on all of \(S_i \otimes_A S_i\). Applying this
  to \(w_{ii}\) and using the two hypotheses,
  \(\mathrm{diag}(w_{ii}) = \delta_i(\kappa_{ii}(w_{ii})) = \delta_i(c_{ii}) = 1\). In
  particular the diagonal hypothesis of
  Theorem~\ref{pa-thm:isDescentCocycle_piAssemblyUnit} is automatic once the components collapse
  onto a cover cocycle.
\end{proof}

\begin{lemma}[The index-wise operators intertwine the two-base structure maps]
  \label{pa-lem:faceA_comp_tensorMap}
  \lean{IsLocalization.AwayCover.faceA₂₃_comp_tensorMap,
    IsLocalization.AwayCover.faceA₁₂_comp_tensorMap,
    IsLocalization.AwayCover.faceA₁₃_comp_tensorMap,
    IsLocalization.AwayCover.componentCollapse_comp_tensorMap}
  \leanok
  \dcref{custom}

  \uses{pa-def:indexwiseMaps, pa-def:componentCollapse, pa-def:tensorAwayAlgebra,
    pa-def:module_descent_cocycle}
  Write \(m \colon B \otimes_A B \to S_i \otimes_A S_j\) for the two-base structure map of
  Definition~\ref{pa-def:tensorAwayAlgebra} (with \(B_1 = B_2 = B\)) and \(m_3 \colon
  B \otimes_A (B \otimes_A B) \to S_i \otimes_A (S_j \otimes_A S_k)\) for its triple-tensor
  analogue. Then the three index-wise faces intertwine the double and triple structure maps
  over the descent cofaces,
  \[
    \mathrm{face}_{\bullet} \circ m = m_3 \circ \partial_{\bullet}
    \qquad (\bullet \in \{12, 13, 23\}) ,
  \]
  and the index-wise collapse intertwines the double structure map with the multiplication:
  \(\kappa_{ij} \circ m = \iota \circ \mu\), where \(\mu \colon B \otimes_A B \to B\) is the
  multiplication and \(\iota \colon B \to T_{ij}\) is the structure map.
\end{lemma}
\begin{proof}
  \leanok
  Each identity is an equality of \(A\)-algebra homomorphisms out of \(B \otimes_A B\)
  (respectively \(B \otimes_A (B \otimes_A B)\)), so it suffices to check it on a pure tensor,
  where both sides are read off from the pure-tensor formulas for the structure maps, faces,
  cofaces, collapse and multiplication. For the \(\{2,3\}\)-face, \(m\) sends
  \(x \otimes y\) to \(\bar x \otimes \bar y\) and \(\mathrm{face}_{23}\) prepends a \(1\),
  giving \(1 \otimes (\bar x \otimes \bar y)\); the same element results from
  \(\partial_{23}(x \otimes y) = 1 \otimes (x \otimes y)\) followed by \(m_3\), using
  \(\bar 1 = 1\). The \(\{1,2\}\)- and \(\{1,3\}\)-faces are identical after tracking which
  tensor slot receives the \(1\). For the collapse, \(\kappa_{ij}(m(x \otimes y)) =
  \iota_1(\bar x)\,\iota_2(\bar y) = \iota(x)\,\iota(y) = \iota(xy) = \iota(\mu(x \otimes y))\),
  since \(\iota_1, \iota_2\) restrict to \(\iota\) on the image of \(B\).
\end{proof}

\section{Degree-zero Amitsur descent and unit production}
\label{pa-sec:amitsurDescentAlgebra}

This section records the two pure-algebra inputs to the faithfully flat descent that closes
the \'etale separatedness of the relative Picard functor: the degree-zero exactness of the
Amitsur complex, in the right-tensor form used to descend a section, and the production of a
unit from the evaluated generator of a trivialized invertible module.

\begin{theorem}[The degree-zero Amitsur equalizer]
  \label{pa-thm:descent_existsUnique_tmul_one}
  \lean{Module.FaithfullyFlat.existsUnique_tmul_one_eq}
  \leanok
  \dcref{custom}

  Let \(A \to B\) be faithfully flat and \(S_0\) an \(A\)-algebra. If an element \(x \in S_0
  \otimes_A B\) has equal images under the two coprojection faces
  \[
    \id_{S_0} \otimes (b \mapsto b \otimes 1), \quad \id_{S_0} \otimes (b \mapsto 1 \otimes b)
    \;\colon\; S_0 \otimes_A B \longrightarrow S_0 \otimes_A (B \otimes_A B),
  \]
  then \(x = s \otimes 1\) for a unique \(s \in S_0\).
\end{theorem}
\begin{proof}
  \leanok
  This is exactness in degree \(0\) of the Amitsur complex of the faithfully flat extension
  \(A \to B\), applied with coefficients in \(S_0\).  Transport \(x\) through the
  commutativity and associativity isomorphisms into the coaction encoding of the tautological
  base-change descent datum on \(S_0\); the two coprojection faces become the two legs of that
  coaction, so the hypothesis says exactly that the coaction difference vanishes on the
  transported element.  Degree-zero Amitsur exactness for a faithfully flat base change then
  produces \(s \in S_0\) whose image \(1 \otimes s\) is the transported element; applying the
  inverse transport gives \(s \otimes 1 = x\).  Uniqueness is faithful flatness: the map \(s
  \mapsto s \otimes 1\), \(S_0 \to S_0 \otimes_A B\), is injective when \(B\) is faithfully
  flat over \(A\).
\end{proof}

\begin{theorem}[Unit production from an evaluated generator]
  \label{pa-thm:isUnit_map_rTensor_generator}
  \lean{Module.isUnit_map_rTensor_generator}
  \leanok
  \dcref{custom}

  Let \(A \to R \to R'\) be a tower with \(R'\) also an \(A\)-algebra and \(\varphi \colon R \to
  R'\) the ring map, carried by an \(A\)-linear avatar.  Let \(M'\) be an \(R\)-module presented
  as trivial by an \(R\)-linear isomorphism \(t \colon M' \cong R\) together with a collapse
  isomorphism \(c \colon M' \cong R \otimes_A M\).  Then, for every surjective \(R'\)-linear map
  \(F \colon R' \otimes_A M \to R'\), the value
  \[
    F\bigl(\varphi_{*}\,c(t^{-1} 1)\bigr) \in R'
  \]
  of \(F\) on the \(\varphi\)-pushforward of the collapsed generator \(c(t^{-1} 1)\) is a unit.
\end{theorem}
\begin{proof}
  \leanok
  Write \(g_0 := c(t^{-1} 1) \in R \otimes_A M\).  Since \(t\) is an isomorphism, every \(x \in
  R \otimes_A M\) is the scalar multiple \(x = t(c^{-1} x) \cdot g_0\); that is, \(g_0\)
  generates \(R \otimes_A M\) as an \(R\)-module.  Pushing forward along \(\varphi\): the map
  \(\varphi_{*} \colon R \otimes_A M \to R' \otimes_A M\) is \(\varphi\)-semilinear, so every
  \(z \in R' \otimes_A M\) is an \(R'\)-multiple of \(g := \varphi_{*} g_0\).  Indeed on a pure
  tensor \(r' \otimes m = r' \cdot (1 \otimes m)\) and \(1 \otimes m = \varphi_{*}(1 \otimes m)
  = \varphi_{*}\bigl(t(c^{-1}(1 \otimes m)) \cdot g_0\bigr) = \varphi\bigl(t(c^{-1}(1 \otimes
  m))\bigr) \cdot g\), and the general element is a finite sum of such.  Now choose, by
  surjectivity of \(F\), an element \(y\) with \(F(y) = 1\), and write \(y = c' \cdot g\) for
  some \(c' \in R'\).  Then \(1 = F(c' \cdot g) = c' \cdot F(g)\), so \(F(g)\) is a unit.
\end{proof}

\begin{definition}[Base change of a descent equivalence]
  \label{pa-def:descentMulEval}
  \lean{Module.descentMulEval}
  \leanok
  \dcref{custom}

  \uses{pa-thm:isUnit_map_rTensor_generator}
  Let \(A \to B' \to R'\) be a tower and \(\mathrm{dE} \colon B' \otimes_A M \cong B'\) a
  \(B'\)-linear descent equivalence.  Its \emph{base change} is the \(R'\)-linear equivalence
  \[
    \mathrm{descentMulEval}(\mathrm{dE}) \colon R' \otimes_A M \;\xrightarrow{\ \sim\ }\; R',
    \qquad r' \otimes m \longmapsto \mathrm{dE}(1 \otimes m) \cdot r',
  \]
  obtained by cancelling the iterated base change \(R' \otimes_{B'} (B' \otimes_A M) \cong R'
  \otimes_A M\), applying \(\mathrm{dE}\) in the second factor, and using the right unit \(R'
  \otimes_{B'} B' \cong R'\).  Being an isomorphism, it is in particular a surjective
  \(R'\)-linear map, the concrete evaluation fed to Theorem~\ref{pa-thm:isUnit_map_rTensor_generator}.
\end{definition}

\section{Additivity of dimension along exact sequences}
\label{pa-sec:finrank_additivity}

The colength computations of the next section, and the Euler characteristic of the
Riemann--Roch ledger, rest on the additivity of the \(k\)-dimension along exact sequences of
finite-dimensional vector spaces. Throughout, \(k\) is a field and all vector spaces are
finite-dimensional over \(k\); \(\dim_k\) is the dimension.

\begin{theorem}[Additivity of dimension along a short exact sequence]
  \label{pa-thm:finrank_add_of_exact}
  \lean{AlgebraicGeometry.finrank_add_of_exact}
  \leanok
  \dcref{custom}

  Let
  \[
    0 \longrightarrow A \xrightarrow{\ f\ } B \xrightarrow{\ g\ } C \longrightarrow 0
  \]
  be a short exact sequence of finite-dimensional \(k\)-vector spaces (\(f\) injective, \(g\)
  surjective, and \(\ker g = \operatorname{im} f\)). Then
  \[
    \dim_k B \;=\; \dim_k A + \dim_k C .
  \]
\end{theorem}
\begin{proof}
  \leanok
  The length of a module is additive along a short exact sequence: for
  \(0 \to A \to B \to C \to 0\) one has \(\ell(B) = \ell(A) + \ell(C)\), where \(\ell\) denotes
  the length of a composition series. For a finite-dimensional vector space over the field
  \(k\), a composition series is a maximal flag of subspaces, so its length equals the
  dimension: \(\ell(V) = \dim_k V\). Substituting this identity for each of \(A\), \(B\),
  \(C\) into the length additivity gives \(\dim_k B = \dim_k A + \dim_k C\).
\end{proof}

\begin{theorem}[Vanishing of the alternating sum of dimensions along a five-term sequence]
  \label{pa-thm:finrank_alt_sum_exact5}
  \lean{AlgebraicGeometry.finrank_alt_sum_eq_zero_of_exact₅}
  \leanok
  \dcref{custom}

  Let
  \[
    0 \longrightarrow V_1 \xrightarrow{\ f_1\ } V_2 \xrightarrow{\ f_2\ } V_3
      \xrightarrow{\ f_3\ } V_4 \xrightarrow{\ f_4\ } V_5 \longrightarrow 0
  \]
  be an exact sequence of finite-dimensional \(k\)-vector spaces (\(f_1\) injective, \(f_4\)
  surjective, and \(\ker f_{i+1} = \operatorname{im} f_i\) for \(i = 1, 2, 3\)). Then the
  alternating sum of dimensions vanishes in \(\Z\):
  \[
    \dim_k V_1 - \dim_k V_2 + \dim_k V_3 - \dim_k V_4 + \dim_k V_5 \;=\; 0 .
  \]
\end{theorem}
\begin{proof}
  \leanok
  Apply the rank--nullity theorem to each of the four maps: for every \(i\),
  \[
    \dim_k V_i \;=\; \dim_k(\ker f_i) + \dim_k(\operatorname{im} f_i) .
  \]
  Injectivity of \(f_1\) gives \(\dim_k(\ker f_1) = 0\); surjectivity of \(f_4\) gives
  \(\dim_k(\operatorname{im} f_4) = \dim_k V_5\); and exactness at \(V_{i+1}\) rewrites the
  kernel as the preceding range, \(\dim_k(\ker f_{i+1}) = \dim_k(\operatorname{im} f_i)\) for
  \(i = 1, 2, 3\). Writing \(r_i = \dim_k(\operatorname{im} f_i)\), these read
  \[
    \dim_k V_1 = r_1,\quad
    \dim_k V_2 = r_1 + r_2,\quad
    \dim_k V_3 = r_2 + r_3,\quad
    \dim_k V_4 = r_3 + \dim_k V_5 .
  \]
  Forming the alternating sum, every \(r_i\) and \(\dim_k V_5\) cancels:
  \[
    \dim_k V_1 - \dim_k V_2 + \dim_k V_3 - \dim_k V_4 + \dim_k V_5
      = r_1 - (r_1 + r_2) + (r_2 + r_3) - (r_3 + \dim_k V_5) + \dim_k V_5 = 0 .
  \]
\end{proof}

\section{The colength of Dedekind quotients}
\label{pa-sec:dedekind_colength}

This section computes the \(k\)-dimension of the quotient of a Dedekind domain \(B\), finitely
generated as a module over a base field \(k\) through a fixed \(k\)-algebra structure, by a
nonzero ideal. The pattern is the classical dévissage: a nonzero principal ideal
\(I = (\pi)\) has all its successive powers separated by copies of the residue quotient
\(B/I\), so \(\dim_k(B/I^e) = e \cdot \dim_k(B/I)\); the Chinese remainder theorem then reduces
an arbitrary ideal to its prime-power factors. Throughout, \(k\) is a field and \(B\) is a
commutative domain with a \(k\)-algebra structure; ideals of \(B\) are written \(I\), \(P\).

\begin{theorem}[One-step filtration jump]
  \label{pa-thm:finrank_quotSucc_jump}
  \lean{AlgebraicGeometry.finrank_quotSucc_jump, AlgebraicGeometry.quotEquivMapPow}
  \leanok
  \dcref{custom}

  Let \(I = (\pi)\) be a nonzero principal ideal of \(B\), with \(B/I\) finite-dimensional over
  \(k\), and let \(n \ge 0\). Multiplication by \(\pi^n\) induces a \(k\)-linear isomorphism
  \[
    B/I \;\xrightarrow{\ \sim\ }\; I^n / I^{n+1} ,
  \]
  so that \(\dim_k\bigl(I^n / I^{n+1}\bigr) = \dim_k(B/I)\).
\end{theorem}
\begin{proof}
  \leanok
  Since \(I = (\pi)\) is principal, \(I^n = (\pi^n)\) and \(I^{n+1} = (\pi^{n+1})\), and
  \(I^{n+1} = I \cdot I^n\) is a \(B\)-submodule of \(I^n\). Multiplication by \(\pi^n\) is a
  \(B\)-linear map \(\mu : B \to I^n\), surjective because every element of \(I^n = (\pi^n)\)
  is a \(B\)-multiple of \(\pi^n\). Composing with the projection \(I^n \twoheadrightarrow
  I^n/I^{n+1}\) yields a surjective \(B\)-linear map \(B \to I^n/I^{n+1}\). Its kernel is
  \(I\): for \(b \in B\), \(\pi^n b \in I^{n+1} = (\pi^{n+1})\) means \(\pi^n b = \pi^{n+1} c\)
  for some \(c\), i.e.\ \(\pi^n(b - \pi c) = 0\); as \(B\) is a domain and \(\pi^n \neq 0\)
  (because \(I \neq 0\)), this forces \(b = \pi c \in (\pi) = I\), and conversely every element
  of \(I\) maps into \(I^{n+1}\). Hence \(\mu\) descends to a \(B\)-linear isomorphism
  \(B/I \cong I^n/I^{n+1}\). Restriction of scalars along \(k \to B\) makes it \(k\)-linear,
  so the two \(k\)-vector spaces have equal dimension.
\end{proof}

\begin{theorem}[Inductive step of the colength formula]
  \label{pa-thm:finrank_quotient_pow_succ}
  \lean{AlgebraicGeometry.finrank_quotient_pow_succ}
  \leanok
  \dcref{custom}

  \uses{pa-thm:finrank_add_of_exact, pa-thm:finrank_quotSucc_jump}
  Let \(I\) be a nonzero principal ideal of \(B\) and \(n \ge 0\), with \(B/I^{n+1}\) and
  \(B/I^n\) finite-dimensional over \(k\). Then
  \[
    \dim_k\bigl(B/I^{n+1}\bigr) \;=\; \dim_k(B/I) + \dim_k\bigl(B/I^n\bigr) .
  \]
\end{theorem}
\begin{proof}
  \leanok
  Because \(I^{n+1} \subseteq I^n\), the identity of \(B\) induces a surjective \(k\)-linear
  projection \(B/I^{n+1} \twoheadrightarrow B/I^n\), whose kernel is the image of \(I^n\) in
  \(B/I^{n+1}\), namely \(I^n/I^{n+1}\). This gives a short exact sequence of
  finite-dimensional \(k\)-vector spaces
  \[
    0 \longrightarrow I^n/I^{n+1} \longrightarrow B/I^{n+1} \longrightarrow B/I^n
      \longrightarrow 0 .
  \]
  By \ref{pa-thm:finrank_add_of_exact}, \(\dim_k(B/I^{n+1}) = \dim_k(I^n/I^{n+1}) +
  \dim_k(B/I^n)\), and by the filtration jump \ref{pa-thm:finrank_quotSucc_jump},
  \(\dim_k(I^n/I^{n+1}) = \dim_k(B/I)\). Combining the two gives the claim.
\end{proof}

\begin{theorem}[Prime-power (DVR) colength formula]
  \label{pa-thm:finrank_quotient_span_pow}
  \lean{AlgebraicGeometry.finrank_quotient_span_pow}
  \leanok
  \dcref{custom}

  \uses{pa-thm:finrank_quotient_pow_succ}
  Let \(I = (\pi)\) be a nonzero principal ideal of \(B\) and \(e \ge 0\), with \(B/I^e\)
  finite-dimensional over \(k\). Then
  \[
    \dim_k\bigl(B/I^e\bigr) \;=\; e \cdot \dim_k(B/I) .
  \]
\end{theorem}
\begin{proof}
  \leanok
  Induction on \(e\). For \(e = 0\), \(I^0\) is the unit ideal \((1) = B\), so \(B/I^0\) is the
  zero module and \(\dim_k(B/I^0) = 0 = 0 \cdot \dim_k(B/I)\). For the step, assume the formula
  for \(e = n\); each quotient \(B/I^n\) is a quotient of \(B/I^{n+1}\), hence
  finite-dimensional, so \ref{pa-thm:finrank_quotient_pow_succ} applies and
  \[
    \dim_k\bigl(B/I^{n+1}\bigr) = \dim_k(B/I) + \dim_k\bigl(B/I^n\bigr)
      = \dim_k(B/I) + n \cdot \dim_k(B/I) = (n+1) \cdot \dim_k(B/I) .
  \]
\end{proof}

For the remaining two results \(B\) is a Dedekind domain, so every nonzero ideal factors
uniquely into height-one primes. For a nonzero ideal \(I\) we write its factorization as
\(I = \prod_P P^{\,\mathrm{ord}_P(I)}\), the product over the finitely many prime factors \(P\)
of \(I\) and \(\mathrm{ord}_P(I) \ge 1\) the multiplicity of \(P\).

\begin{definition}[The Chinese remainder decomposition of a Dedekind quotient]
  \label{pa-def:quotAlgEquivPiFactors}
  \lean{AlgebraicGeometry.quotAlgEquivPiFactors}
  \leanok
  \dcref{custom}

  Let \(B\) be a Dedekind domain, a \(k\)-algebra, and \(I \neq 0\) an ideal with prime
  factorization \(I = \prod_{P \mid I} P^{\,\mathrm{ord}_P(I)}\). The canonical map
  \[
    B/I \;\xrightarrow{\ \sim\ }\; \prod_{P \mid I} B/P^{\,\mathrm{ord}_P(I)}
  \]
  is an isomorphism of \(k\)-algebras.
\end{definition}
\begin{proof}
  \leanok
  Distinct prime factors \(P \neq P'\) of \(I\) are nonzero primes of a Dedekind domain,
  hence maximal, hence comaximal, and powers of comaximal ideals remain comaximal; the
  Chinese remainder theorem therefore identifies \(B/\prod_P P^{\,\mathrm{ord}_P(I)}\)
  with the product of the \(B/P^{\,\mathrm{ord}_P(I)}\), by the canonical map. Each
  component of the canonical map is the quotient map
  \(B/I \to B/P^{\,\mathrm{ord}_P(I)}\), a \(k\)-algebra homomorphism, so the
  isomorphism is \(k\)-algebraic.
\end{proof}

\begin{theorem}[Multi-prime colength via the Chinese remainder theorem]
  \label{pa-thm:finrank_quotient_eq_sum_factors_pow}
  \lean{AlgebraicGeometry.finrank_quotient_eq_sum_factors_pow}
  \leanok
  \dcref{custom}

  Let \(B\) be a Dedekind domain, a \(k\)-algebra, and let \(I \neq 0\) be an ideal with each
  prime-power quotient \(B/P^{\,\mathrm{ord}_P(I)}\) finite-dimensional over \(k\). Then
  \[
    \dim_k(B/I) \;=\; \sum_{P \mid I} \dim_k\bigl(B/P^{\,\mathrm{ord}_P(I)}\bigr) ,
  \]
  the sum ranging over the prime factors \(P\) of \(I\).
\end{theorem}
\begin{proof}
  \leanok
  \uses{pa-def:quotAlgEquivPiFactors}
  The Chinese remainder decomposition \ref{pa-def:quotAlgEquivPiFactors} identifies
  \(B/I\) with \(\prod_{P \mid I} B/P^{\,\mathrm{ord}_P(I)}\) as \(k\)-vector spaces.
  The dimension of a finite product of vector spaces is the sum of the dimensions of
  the factors, which is the stated formula.
\end{proof}

\begin{theorem}[Dedekind colength formula, ord-weighted form]
  \label{pa-thm:finrank_quotient_span_eq_sum_ord}
  \lean{AlgebraicGeometry.finrank_quotient_span_eq_sum_ord}
  \leanok
  \dcref{custom}

  \uses{pa-thm:finrank_quotient_eq_sum_factors_pow, pa-thm:finrank_quotient_span_pow}
  Let \(B\) be a Dedekind domain, a \(k\)-algebra, and let \(f \neq 0\) be an element whose
  height-one prime factors \(P\) of the principal ideal \((f)\) are all principal (for instance
  when \(B\) is a principal ideal domain), with each prime-power quotient finite-dimensional
  over \(k\). Then
  \[
    \dim_k\bigl(B/(f)\bigr) \;=\; \sum_{P \mid (f)} \mathrm{ord}_P(f) \cdot \dim_k(B/P) ,
  \]
  where \(\mathrm{ord}_P(f) = \mathrm{ord}_P\bigl((f)\bigr)\) is the multiplicity of \(P\) in
  the factorization of \((f)\).
\end{theorem}
\begin{proof}
  \leanok
  Since \(f \neq 0\), the ideal \((f)\) is nonzero, and \ref{pa-thm:finrank_quotient_eq_sum_factors_pow}
  gives \(\dim_k(B/(f)) = \sum_{P \mid (f)} \dim_k(B/P^{\,\mathrm{ord}_P(f)})\). Each prime
  factor \(P\) is a nonzero principal ideal by hypothesis, so the prime-power formula
  \ref{pa-thm:finrank_quotient_span_pow} applies to it and yields
  \(\dim_k(B/P^{\,\mathrm{ord}_P(f)}) = \mathrm{ord}_P(f) \cdot \dim_k(B/P)\). Substituting into
  the sum gives the claim.
\end{proof}

\section{The Dedekind colength through localization at a prime}
\label{pa-sec:localized_colength}

The ord-weighted formula \ref{pa-thm:finrank_quotient_span_eq_sum_ord} assumes each prime
factor of \((f)\) is principal --- true in a principal ideal domain, but false in a
general Dedekind domain. This section removes the principality hypothesis: each
prime-power block \(B/P^e\) is transported to the localization of \(B\) at \(P\), a
discrete valuation ring, where the maximal ideal \emph{is} principal and the prime-power
formula \ref{pa-thm:finrank_quotient_span_pow} applies. The bookkeeping needed for the
transport --- that multiplicities of prime factors are computed by membership in prime
powers, and are preserved by the localization --- is proved along the way. Throughout,
\(k\) is a field; \(B\) is a domain with a \(k\)-algebra structure, Dedekind where
stated; and for a nonzero ideal \(I \subseteq B\) of a Dedekind domain,
\(\mathrm{ord}_P(I)\) denotes the multiplicity of the prime \(P\) in the factorization of
\(I\), with \(\mathrm{ord}_P(f) = \mathrm{ord}_P\bigl((f)\bigr)\) for an element
\(f \neq 0\).

\begin{theorem}[Finiteness d\'evissage along principal prime powers]
  \label{pa-thm:moduleFinite_quotient_pow_of_isPrincipal}
  \lean{AlgebraicGeometry.moduleFinite_quotient_pow_of_isPrincipal}
  \leanok
  \dcref{custom}

  Let \(B\) be a domain over \(k\) and \(I = (\pi)\) a nonzero principal ideal with
  \(B/I\) finite-dimensional over \(k\). Then \(B/I^n\) is finite-dimensional over
  \(k\) for every \(n \ge 0\).
\end{theorem}
\begin{proof}
  \leanok
  \uses{pa-thm:finrank_quotSucc_jump}
  Induction on \(n\). For \(n = 0\) the quotient \(B/I^0 = B/B\) is zero. For the step,
  the projection \(B/I^{n+1} \twoheadrightarrow B/I^n\) is surjective with kernel
  \(I^n/I^{n+1}\), which is \(k\)-isomorphic to \(B/I\) by the filtration jump
  \ref{pa-thm:finrank_quotSucc_jump}, hence finite-dimensional. Thus \(B/I^{n+1}\) is an
  extension of the finite-dimensional \(B/I^n\) (inductive hypothesis) by the
  finite-dimensional \(I^n/I^{n+1}\), and is finite-dimensional.
\end{proof}

\begin{lemma}[The adic valuation counts the multiplicity]
  \label{pa-lem:adic_valuation_count}
  \lean{AlgebraicGeometry.intValuation_eq_exp_neg_count}
  \leanok
  \dcref{custom}

  Let \(B\) be a Dedekind domain, \(P\) a nonzero prime of \(B\), and \(f \neq 0\).
  Then the \(P\)-adic valuation of \(f\) equals the multiplicity of \(P\) in the prime
  factorization of the principal ideal \((f)\):
  \[
    v_P(f) \;=\; \mathrm{ord}_P(f) .
  \]
\end{lemma}
\begin{proof}
  \leanok
  In a Dedekind domain every nonzero ideal factors uniquely into nonzero primes, and
  the \(P\)-adic valuation of a nonzero element \(f\) is by definition the exponent of
  \(P\) in the factorization of \((f)\). Uniqueness of the factorization identifies
  that exponent with the number of occurrences of \(P\) in the multiset of prime
  factors of \((f)\), which is \(\mathrm{ord}_P(f)\).
\end{proof}

\begin{lemma}[Membership in prime powers reads off the multiplicity]
  \label{pa-lem:mem_pow_iff_le_count}
  \lean{AlgebraicGeometry.mem_pow_iff_le_count}
  \leanok
  \dcref{custom}

  \uses{pa-lem:adic_valuation_count}
  Let \(B\) be a Dedekind domain, \(P\) a nonzero prime, \(f \neq 0\) and \(n \ge 0\).
  Then
  \[
    f \in P^n \quad\Longleftrightarrow\quad n \le \mathrm{ord}_P(f) .
  \]
\end{lemma}
\begin{proof}
  \leanok
  Membership \(f \in P^n\) means \((f) \subseteq P^n\), and in a Dedekind domain to
  contain is to divide, so this holds exactly when \(v_P(f) \ge n\) --- the
  valuation-theoretic characterization of the powers of \(P\). By
  \ref{pa-lem:adic_valuation_count}, \(v_P(f) = \mathrm{ord}_P(f)\), giving the claim.
\end{proof}

\begin{theorem}[Multiplicities are preserved by localization at the prime]
  \label{pa-thm:count_factors_span_localization}
  \lean{AlgebraicGeometry.count_factors_span_algebraMap}
  \leanok
  \dcref{custom}

  \uses{pa-lem:mem_pow_iff_le_count}
  Let \(B\) be a Dedekind domain, \(P\) a nonzero prime of \(B\), and \(S\) a Dedekind
  local ring which is a localization of \(B\) at \(P\) --- for instance the
  localization \(B_P\) itself, a discrete valuation ring. Write \(\mathfrak m\) for the
  maximal ideal of \(S\). Then for every \(f \neq 0\) in \(B\),
  \[
    \mathrm{ord}_{\mathfrak m}\bigl(f \cdot S\bigr) \;=\; \mathrm{ord}_P(f) :
  \]
  the multiplicity of \(\mathfrak m\) in the factorization of the extended principal
  ideal equals the multiplicity of \(P\) in the factorization of \((f)\).
\end{theorem}
\begin{proof}
  \leanok
  The localization map \(B \to S\) is injective (its kernel consists of elements killed
  by some \(s \notin P\), and \(B\) is a domain), so the image of \(f\) in \(S\) is
  nonzero and both multiplicities are defined. Under the localization, the powers of
  \(\mathfrak m\) contract to the powers of \(P\): \(\mathfrak m^n \cap B = P^n\)
  (the \(P\)-primary ideal \(P^n\) is saturated with respect to the complement of
  \(P\)). Hence for every \(n \ge 0\),
  \[
    f \cdot S \subseteq \mathfrak m^n
    \quad\Longleftrightarrow\quad
    f \in P^n .
  \]
  By \ref{pa-lem:mem_pow_iff_le_count}, applied in \(S\) (a Dedekind domain, with
  \(\mathfrak m\) its nonzero prime) and in \(B\), the left side holds exactly when
  \(n \le \mathrm{ord}_{\mathfrak m}(f \cdot S)\) and the right side exactly when
  \(n \le \mathrm{ord}_P(f)\). Two natural numbers with the same lower sets are equal.
\end{proof}

\begin{definition}[The localization isomorphisms, \(k\)-linearly]
  \label{pa-def:quotPowLinearEquiv}
  \lean{AlgebraicGeometry.quotPowLinearEquiv, AlgebraicGeometry.residueQuotLinearEquiv}
  \leanok
  \dcref{custom}

  Let \(B\) be a \(k\)-algebra, \(P\) a maximal ideal of \(B\), and \(S\) a local
  \(k\)-algebra which is a localization of \(B\) at \(P\) compatibly with the
  \(k\)-structures, with maximal ideal \(\mathfrak m\). For every \(n \ge 0\) the
  canonical ring isomorphism
  \[
    B/P^n \;\xrightarrow{\ \sim\ }\; S/\mathfrak m^n ,
  \]
  induced by the localization map, is \(k\)-linear; in particular (\(n = 1\)) the
  residue isomorphism \(B/P \cong S/\mathfrak m\) is \(k\)-linear.
\end{definition}
\begin{proof}
  \leanok
  The composite \(B \to S \to S/\mathfrak m^n\) kills \(P^n\), and the induced map
  \(B/P^n \to S/\mathfrak m^n\) is bijective: every element of the complement of \(P\)
  acts invertibly on \(B/P^n\) (the ideal \(P^n\) is \(P\)-primary with \(P\) maximal,
  so \(B/P^n\) is local with maximal ideal \(P/P^n\), and elements outside \(P\) map to
  units), so the localization map factors through an inverse. The maps are
  \(k\)-algebra maps because the \(k\)-structure of \(S\) is that of \(B\) composed
  with the localization; hence the isomorphism is \(k\)-linear.
\end{proof}

\begin{theorem}[Per-prime finiteness through the localization]
  \label{pa-thm:moduleFinite_quotient_pow_atPrime}
  \lean{AlgebraicGeometry.moduleFinite_quotient_pow_atPrime}
  \leanok
  \dcref{custom}

  \uses{pa-def:quotPowLinearEquiv}
  Let \(B\) be a Dedekind domain over \(k\), \(P\) a nonzero maximal ideal with
  residue \(B/P\) finite-dimensional over \(k\), and \(S\) a domain over \(k\) which
  is a localization of \(B\) at \(P\) compatibly with the \(k\)-structures. Then
  \(B/P^n\) is finite-dimensional over \(k\) for every \(n \ge 0\).
\end{theorem}
\begin{proof}
  \leanok
  \uses{pa-thm:moduleFinite_quotient_pow_of_isPrincipal}
  The localization \(S\) of the Dedekind domain \(B\) at the nonzero prime \(P\) is a
  discrete valuation ring; its maximal ideal \(\mathfrak m\) is nonzero and principal.
  By \ref{pa-def:quotPowLinearEquiv} the residue \(S/\mathfrak m \cong B/P\) is
  finite-dimensional, so the d\'evissage
  \ref{pa-thm:moduleFinite_quotient_pow_of_isPrincipal}, applied in \(S\) to the nonzero
  principal ideal \(\mathfrak m\), makes every \(S/\mathfrak m^n\) finite-dimensional;
  transporting back along \ref{pa-def:quotPowLinearEquiv} finishes.
\end{proof}

\begin{theorem}[The per-prime colength collapse]
  \label{pa-thm:finrank_quotient_pow_atPrime}
  \lean{AlgebraicGeometry.finrank_quotient_pow_atPrime}
  \leanok
  \dcref{custom}

  \uses{pa-def:quotPowLinearEquiv}
  In the situation of \ref{pa-thm:moduleFinite_quotient_pow_atPrime},
  \[
    \dim_k\bigl(B/P^n\bigr) \;=\; n \cdot \dim_k(B/P)
    \qquad (n \ge 0).
  \]
\end{theorem}
\begin{proof}
  \leanok
  \uses{pa-thm:moduleFinite_quotient_pow_of_isPrincipal, pa-thm:finrank_quotient_span_pow}
  As in \ref{pa-thm:moduleFinite_quotient_pow_atPrime}, \(S\) is a discrete valuation
  ring with nonzero principal maximal ideal \(\mathfrak m\), the quotients
  \(S/\mathfrak m^n\) are finite-dimensional
  (\ref{pa-thm:moduleFinite_quotient_pow_of_isPrincipal}), and the prime-power formula
  \ref{pa-thm:finrank_quotient_span_pow} in \(S\) gives
  \(\dim_k(S/\mathfrak m^n) = n \cdot \dim_k(S/\mathfrak m)\). The \(k\)-linear
  identifications \(B/P^n \cong S/\mathfrak m^n\) and \(B/P \cong S/\mathfrak m\)
  (\ref{pa-def:quotPowLinearEquiv}) transport this to the claim.
\end{proof}

\begin{lemma}[Prime factors of a principal ideal]
  \label{pa-lem:prime_factor_span}
  \lean{AlgebraicGeometry.prime_factor_span}
  \leanok
  \dcref{custom}

  Let \(B\) be a Dedekind domain and \(f \in B\). Every prime factor \(P\) of the
  ideal \((f)\) is a nonzero prime ideal of \(B\) containing \(f\).
\end{lemma}
\begin{proof}
  \leanok
  A factor occurring in the prime factorization of \((f)\) is by definition a nonzero
  prime ideal, and it divides \((f)\); in a Dedekind domain to divide is to contain,
  so \((f) \subseteq P\), i.e.\ \(f \in P\).
\end{proof}

\begin{theorem}[Finiteness of the Dedekind colength module]
  \label{pa-thm:moduleFinite_quotient_span}
  \lean{AlgebraicGeometry.moduleFinite_quotient_span}
  \leanok
  \dcref{custom}

  \uses{pa-def:quotAlgEquivPiFactors}
  Let \(B\) be a Dedekind domain over \(k\) and \(f \neq 0\) an element such that the
  residue \(B/P\) is finite-dimensional over \(k\) for every prime factor \(P\) of
  \((f)\). Then \(B/(f)\) is finite-dimensional over \(k\).
\end{theorem}
\begin{proof}
  \leanok
  \uses{pa-lem:prime_factor_span, pa-thm:moduleFinite_quotient_pow_atPrime}
  The ideal \((f)\) is nonzero, so the Chinese remainder decomposition
  \ref{pa-def:quotAlgEquivPiFactors} identifies \(B/(f)\) with the finite product
  \(\prod_{P \mid (f)} B/P^{\,\mathrm{ord}_P(f)}\). Each factor is
  finite-dimensional: \(P\) is a nonzero prime (\ref{pa-lem:prime_factor_span}), hence
  maximal, with finite-dimensional residue by hypothesis, so
  \ref{pa-thm:moduleFinite_quotient_pow_atPrime} applies with \(S = B_P\) (a domain,
  being a subring of the fraction field of \(B\)). A finite product of
  finite-dimensional spaces is finite-dimensional.
\end{proof}

\begin{theorem}[Dedekind colength formula, localized form]
  \label{pa-thm:finrank_quotient_span_eq_sum_count}
  \lean{AlgebraicGeometry.finrank_quotient_span_eq_sum_count}
  \leanok
  \dcref{custom}

  \uses{pa-thm:finrank_quotient_eq_sum_factors_pow}
  Let \(B\) be a Dedekind domain over \(k\) and \(f \neq 0\) an element such that the
  residue \(B/P\) is finite-dimensional over \(k\) for every prime factor \(P\) of
  \((f)\). Then
  \[
    \dim_k\bigl(B/(f)\bigr) \;=\;
      \sum_{P \mid (f)} \mathrm{ord}_P(f) \cdot \dim_k(B/P) ,
  \]
  the sum ranging over the prime factors of \((f)\). No principality of the primes
  \(P\) is assumed: this is the generalization of
  \ref{pa-thm:finrank_quotient_span_eq_sum_ord} to arbitrary Dedekind domains.
\end{theorem}
\begin{proof}
  \leanok
  \uses{pa-lem:prime_factor_span, pa-thm:moduleFinite_quotient_pow_atPrime,
    pa-thm:finrank_quotient_pow_atPrime}
  Each prime factor \(P\) of \((f)\) is nonzero (\ref{pa-lem:prime_factor_span}), hence
  maximal, and \ref{pa-thm:moduleFinite_quotient_pow_atPrime} with \(S = B_P\) makes each
  prime-power quotient \(B/P^{\,\mathrm{ord}_P(f)}\) finite-dimensional. The
  multi-prime formula \ref{pa-thm:finrank_quotient_eq_sum_factors_pow} then gives
  \(\dim_k(B/(f)) = \sum_{P \mid (f)} \dim_k\bigl(B/P^{\,\mathrm{ord}_P(f)}\bigr)\),
  and the per-prime collapse \ref{pa-thm:finrank_quotient_pow_atPrime} (again through
  \(S = B_P\)) evaluates each summand as
  \(\mathrm{ord}_P(f) \cdot \dim_k(B/P)\).
\end{proof}

\section{The diagonal ideal of an \'etale-coordinatised algebra}
\label{pa-sec:diagonal_ideal}

Let \(k\) be a commutative ring and \(B\) a commutative \(k\)-algebra carrying a compatible
structure of algebra over the polynomial ring \(k[X]\), that is, a tower of algebra maps
\(k \to k[X] \to B\). Write \(u := \) the image of \(X\) in \(B\) for the \emph{\'etale
coordinate}. The two tensor squares \(B \otimes_k B\) and \(B \otimes_{k[X]} B\) carry the two
multiplication maps \(\mu_k \colon B \otimes_k B \to B\) and \(\mu_{k[X]} \colon B
\otimes_{k[X]} B \to B\), \(x \otimes y \mapsto xy\); their kernels are the \emph{\(k\)-diagonal
ideal} and the \emph{\(k[X]\)-diagonal ideal}. The geometric picture is that \(\Spec B\) is a
standard-smooth chart of a relative curve, \(u\) cuts it out over \(\mathbb A^1 = \Spec k[X]\),
and on a basic open containing the whole diagonal of \(\Spec B \times_k \Spec B\) that diagonal
is the vanishing locus of the single equation \(u \otimes 1 - 1 \otimes u\). This section is the
pure commutative-algebra core computing the \(k\)-diagonal ideal after localising away from an
idempotent lift.

\begin{definition}[The diagonal equation and the base-change collapse]
  \label{pa-def:diagonalSetup}
  \lean{AlgebraicJacobian.Diagonal.coord, AlgebraicJacobian.Diagonal.diagGen,
    AlgebraicJacobian.Diagonal.baseChange, AlgebraicJacobian.Diagonal.lmul'_diagGen,
    AlgebraicJacobian.Diagonal.baseChange_diagGen, AlgebraicJacobian.Diagonal.lmul'_baseChange}
  \leanok
  \dcref{custom}

  The \emph{diagonal equation} is
  \[
    \delta \;:=\; u \otimes 1 - 1 \otimes u \;\in\; B \otimes_k B .
  \]
  The \emph{base-change collapse} is the \(k\)-algebra homomorphism
  \[
    \pi \colon B \otimes_k B \longrightarrow B \otimes_{k[X]} B, \qquad x \otimes y \mapsto x
    \otimes y,
  \]
  the surjection identifying the two tensor squares. It satisfies:
  \begin{itemize}
    \item \(\mu_k(\delta) = 0\): the diagonal equation lies in the \(k\)-diagonal ideal;
    \item \(\pi(\delta) = 0\): the collapse kills the diagonal equation, because \(u \otimes 1 =
      1 \otimes u\) already in \(B \otimes_{k[X]} B\) (here \(u = X \cdot 1\) is a scalar of the
      tensor product);
    \item \(\mu_{k[X]} \circ \pi = \mu_k\): the two multiplications agree through the collapse.
  \end{itemize}
\end{definition}
\begin{proof}
  \leanok
  The three identities are checked on pure tensors. For \(\mu_k(\delta)\): \(\mu_k(u \otimes 1) -
  \mu_k(1 \otimes u) = u - u = 0\). For \(\pi(\delta)\): \(u \otimes 1 = (X \cdot 1) \otimes 1 = 1
  \otimes (X \cdot 1) = 1 \otimes u\) in \(B \otimes_{k[X]} B\), since \(X\) is a scalar of the
  \(k[X]\)-tensor product and passes across the tensor. For the factorisation: both sides send
  \(x \otimes y\) to \(xy\), and \(\pi\), \(\mu_{k[X]}\), \(\mu_k\) are additive, so agreement on
  pure tensors suffices.
\end{proof}

\subsection{The kernel of the collapse (deliverable (a))}

\begin{lemma}[Polynomial telescoping]
  \label{pa-lem:diagonal_telescoping}
  \lean{AlgebraicJacobian.Diagonal.algebraMap_tmul_sub_mem_span}
  \leanok
  \dcref{custom}

  \uses{pa-def:diagonalSetup}
  Write \(g \colon k[X] \to B\) for the coordinate map, so \(g(p) = p(u)\). For every polynomial
  \(p \in k[X]\),
  \[
    g(p) \otimes 1 - 1 \otimes g(p) \;\in\; (\delta),
  \]
  the ideal generated by the diagonal equation.
\end{lemma}
\begin{proof}
  \leanok
  Since \(g(p) = p(u)\) is the evaluation of \(p\) at \(u\), the two tensor factors are
  themselves evaluations: \(g(p) \otimes 1 = p(u \otimes 1)\) and \(1 \otimes g(p) = p(1 \otimes
  u)\), each obtained by applying the algebra map \(a \mapsto p(a)\) to the element \(u \otimes 1\)
  respectively \(1 \otimes u\) of \(B \otimes_k B\) (using that the two inclusions \(B \to B
  \otimes_k B\) are algebra maps). The difference of the evaluations of a polynomial at two
  points \(a, b\) is divisible by \(a - b\); here \(a - b = (u \otimes 1) - (1 \otimes u) =
  \delta\), so \(p(u \otimes 1) - p(1 \otimes u) \in (\delta)\).
\end{proof}

\begin{lemma}[The section of the collapse]
  \label{pa-lem:diagonal_section}
  \lean{AlgebraicJacobian.Diagonal.sectionMap, AlgebraicJacobian.Diagonal.qRight,
    AlgebraicJacobian.Diagonal.sectionMap_baseChange}
  \leanok
  \dcref{custom}

  \uses{pa-def:diagonalSetup, pa-lem:diagonal_telescoping}
  There is a \(k[X]\)-algebra homomorphism
  \[
    \psi \colon B \otimes_{k[X]} B \longrightarrow (B \otimes_k B)/(\delta), \qquad x \otimes y
    \mapsto [x \otimes y],
  \]
  and it is a section of the collapse in the sense that \(\psi \circ \pi\) is the quotient map
  \(B \otimes_k B \to (B \otimes_k B)/(\delta)\).
\end{lemma}
\begin{proof}
  \leanok
  The right inclusion \(b \mapsto [1 \otimes b]\) is a ring homomorphism \(B \to (B \otimes_k
  B)/(\delta)\); it is moreover a \(k[X]\)-algebra map, because the two ways of scaling by \(p \in
  k[X]\) — through \([1 \otimes g(p)]\) and through \([g(p) \otimes 1]\) (the ambient
  \(k[X]\)-structure) — agree in the quotient by the telescoping
  Lemma~\ref{pa-lem:diagonal_telescoping}. Together with the left inclusion \(x \mapsto [x \otimes
  1]\), which commutes with it, the universal property of \(B \otimes_{k[X]} B\) produces the
  \(k[X]\)-algebra map \(\psi\) with \(\psi(x \otimes y) = [x \otimes y]\). Both \(\psi \circ \pi\)
  and the quotient map send a pure tensor \(x \otimes y\) to \([x \otimes y]\), and both are
  additive, so they coincide.
\end{proof}

\begin{theorem}[The kernel of the base-change collapse]
  \label{pa-thm:ker_baseChange}
  \lean{AlgebraicJacobian.Diagonal.ker_baseChange}
  \leanok
  \dcref{custom}

  \uses{pa-def:diagonalSetup, pa-lem:diagonal_section}
  The kernel of the collapse \(\pi \colon B \otimes_k B \to B \otimes_{k[X]} B\) is exactly the
  principal ideal generated by the diagonal equation:
  \[
    \ker \pi \;=\; (\delta) \;=\; (u \otimes 1 - 1 \otimes u).
  \]
\end{theorem}
\begin{proof}
  \leanok
  The inclusion \((\delta) \subseteq \ker \pi\) is \(\pi(\delta) = 0\)
  (Definition~\ref{pa-def:diagonalSetup}). Conversely, if \(\pi(z) = 0\) then, applying the section
  \(\psi\) of Lemma~\ref{pa-lem:diagonal_section}, the class of \(z\) in the quotient is \([z] =
  \psi(\pi(z)) = \psi(0) = 0\), so \(z \in (\delta)\).
\end{proof}

\subsection{The idempotent generator and the localised principality (deliverables (b), (c))}

\begin{theorem}[Idempotent generator of the {\(k[X]\)}-diagonal ideal]
  \label{pa-thm:exists_idempotent_lift}
  \lean{AlgebraicJacobian.Diagonal.exists_idempotent_lift}
  \leanok
  \dcref{custom}

  \uses{pa-def:diagonalSetup}
  Suppose \(k[X] \to B\) is formally unramified and essentially of finite type. Then the
  \(k[X]\)-diagonal ideal \(\ker \mu_{k[X]}\) is generated by an idempotent element, and this
  idempotent is \(\pi(\mathrm{elift})\) for some lift \(\mathrm{elift} \in B \otimes_k B\):
  \[
    \ker \mu_{k[X]} \;=\; (\pi\,\mathrm{elift}), \qquad (\pi\,\mathrm{elift})^2 = \pi\,\mathrm{elift}.
  \]
\end{theorem}
\begin{proof}
  \leanok
  The \(k[X]\)-diagonal ideal \(I' := \ker \mu_{k[X]}\) has \(I'/I'^2 \cong \Omega_{B/k[X]}\), the
  module of K\"ahler differentials, which vanishes because \(k[X] \to B\) is formally unramified;
  hence \(I'\) is idempotent, \(I' = I'^2\). Essential finite type makes \(I'\) finitely
  generated. A finitely generated idempotent ideal is generated by an idempotent element \(e\)
  (\(I' = (e)\) with \(e^2 = e\)). Finally \(\pi\) is surjective, so \(e = \pi(\mathrm{elift})\)
  for some \(\mathrm{elift} \in B \otimes_k B\).
\end{proof}

\begin{theorem}[Localised principality of the \(k\)-diagonal ideal]
  \label{pa-thm:ker_map_localization_eq}
  \lean{AlgebraicJacobian.Diagonal.ker_map_localization_eq}
  \leanok
  \dcref{custom}

  \uses{pa-def:diagonalSetup, pa-thm:ker_baseChange, pa-thm:exists_idempotent_lift}
  Let \(\mathrm{elift} \in B \otimes_k B\) be a lift with \(\pi(\mathrm{elift})\) idempotent and
  \(\ker \mu_{k[X]} = (\pi\,\mathrm{elift})\). Then for any localisation \(L\) of \(B \otimes_k B\)
  away from \(1 - \mathrm{elift}\), the \(k\)-diagonal ideal, extended to \(L\), is the principal
  ideal generated by the diagonal equation:
  \[
    (\ker \mu_k)\cdot L \;=\; (\delta)\cdot L .
  \]
\end{theorem}
\begin{proof}
  \leanok
  First, in \(L\) the lift is a multiple of \(\delta\). Indeed \(\mathrm{elift}\,(1 -
  \mathrm{elift}) \in \ker \pi = (\delta)\) (by Theorem~\ref{pa-thm:ker_baseChange}, since
  \(\pi(\mathrm{elift}(1-\mathrm{elift})) = e(1-e) = 0\) as \(e = \pi(\mathrm{elift})\) is
  idempotent), while \(1 - \mathrm{elift}\) is a unit of \(L\); dividing shows \(\mathrm{elift} \in
  (\delta)\cdot L\).

  Now \((\delta)\cdot L \subseteq (\ker \mu_k)\cdot L\) because \(\delta \in \ker \mu_k\)
  (\(\mu_k(\delta) = 0\)). For the reverse inclusion it suffices, by \(\ker \mu_k \subseteq
  (\delta) + (\mathrm{elift})\), that both \((\delta)\) and \((\mathrm{elift})\) map into
  \((\delta)\cdot L\); the first is clear and the second is the previous paragraph. The
  containment \(\ker \mu_k \subseteq (\delta) + (\mathrm{elift})\) holds because for \(x\) with
  \(\mu_k(x) = 0\) one has \(\pi(x) \in \ker \mu_{k[X]} = (\pi\,\mathrm{elift})\) (using
  \(\mu_{k[X]} \circ \pi = \mu_k\)), so \(\pi(x) = \pi(\mathrm{elift}\cdot w)\) for some \(w\);
  then \(x - \mathrm{elift}\cdot w \in \ker \pi = (\delta)\) and \(x = (x - \mathrm{elift}\,w) +
  \mathrm{elift}\,w \in (\delta) + (\mathrm{elift})\).
\end{proof}

\begin{theorem}[Packaged diagonal localisation]
  \label{pa-thm:exists_diagonalLocalization}
  \lean{AlgebraicJacobian.Diagonal.exists_diagonalLocalization,
    AlgebraicJacobian.Diagonal.lmul'_lift}
  \leanok
  \dcref{custom}

  \uses{pa-thm:ker_baseChange, pa-thm:exists_idempotent_lift, pa-thm:ker_map_localization_eq}
  Suppose \(k[X] \to B\) is formally unramified and essentially of finite type. Then there is a
  lift \(\mathrm{elift} \in B \otimes_k B\) with \(\mu_k(\mathrm{elift}) = 0\), with
  \(\pi(\mathrm{elift})\) idempotent generating \(\ker \mu_{k[X]}\), and with the localised
  principality of Theorem~\ref{pa-thm:ker_map_localization_eq} holding on the basic open away from
  \(1 - \mathrm{elift}\).
\end{theorem}
\begin{proof}
  \leanok
  Take \(\mathrm{elift}\) from Theorem~\ref{pa-thm:exists_idempotent_lift}. Then
  \(\mu_k(\mathrm{elift}) = \mu_{k[X]}(\pi\,\mathrm{elift}) = 0\), since \(\pi(\mathrm{elift}) \in
  \ker \mu_{k[X]}\) and \(\mu_{k[X]} \circ \pi = \mu_k\) (Definition~\ref{pa-def:diagonalSetup}); the
  idempotent-generation is Theorem~\ref{pa-thm:exists_idempotent_lift}, and the localised
  principality is Theorem~\ref{pa-thm:ker_map_localization_eq}.
\end{proof}

\section{The diagonal-equation regularity engine}
\label{pa-sec:diagonal_regular}

The regularity input needed alongside the diagonal ideal is the fact that the diagonal
equation is a nonzerodivisor, and this holds uniformly in an arbitrary \emph{second} tensor
factor, with no integrality, fibre, or transcendence hypothesis — so it survives at every
point of a non-integral product.

\begin{lemma}[A flat ring homomorphism preserves nonzerodivisors]
  \label{pa-lem:flat_mem_nonZeroDivisors}
  \lean{RingHom.Flat.mem_nonZeroDivisors}
  \leanok
  \dcref{custom}

  Let \(f \colon R \to S\) be a flat homomorphism of commutative rings. If \(r \in R\) is a
  nonzerodivisor, then \(f(r) \in S\) is a nonzerodivisor.
\end{lemma}
\begin{proof}
  \leanok
  View \(S\) as an \(R\)-module through \(f\); flatness makes it a flat \(R\)-module.
  Multiplication by a nonzerodivisor \(r\) is injective on a flat module, i.e.\ \(r\) acts
  regularly on \(S\). Since \(r \cdot s = f(r)\, s\) for \(s \in S\), the map \(s \mapsto f(r)\, s\)
  is injective, and by commutativity \(f(r)\) is a nonzerodivisor of \(S\).
\end{proof}

\begin{theorem}[The regularity engine]
  \label{pa-thm:diagonal_regularity_engine}
  \lean{AlgebraicJacobian.Diagonal.tmul_one_sub_one_tmul_mem_nonZeroDivisors,
    AlgebraicJacobian.Diagonal.diagGen_mem_nonZeroDivisors}
  \leanok
  \dcref{custom}

  \uses{pa-def:diagonalSetup, pa-lem:flat_mem_nonZeroDivisors}
  Assume in addition that \(B\) is flat over \(k[X]\). Then for \emph{any} \(k\)-algebra \(A\) and
  \emph{any} \(b \in A\), the element
  \[
    u \otimes 1 - 1 \otimes b \;\in\; B \otimes_k A
  \]
  is a nonzerodivisor. In particular (taking \(A = B\), \(b = u\)) the diagonal equation \(u
  \otimes 1 - 1 \otimes u\) is a nonzerodivisor of \(B \otimes_k B\).
\end{theorem}
\begin{proof}
  \leanok
  The linear polynomial \(X - b \in A[X]\) is monic, hence a nonzerodivisor of \(A[X]\). Consider
  the composite
  \[
    \Theta \colon A[X] \;\xrightarrow{\ \sim\ }\; A \otimes_k k[X]
    \;\xrightarrow{\ \mathrm{id} \otimes (k[X] \to B)\ }\; A \otimes_k B
    \;\xrightarrow{\ \sim\ }\; B \otimes_k A,
  \]
  which sends \(X - b\) to \(u \otimes 1 - 1 \otimes b\). Its middle arrow is the base change
  along \(A\) of the flat coordinate map \(k[X] \to B\), hence flat; the two outer arrows are
  algebra isomorphisms (the polynomial-tensor identification and the commutativity of the tensor
  product), hence flat; so \(\Theta\) is a flat ring homomorphism. By
  Lemma~\ref{pa-lem:flat_mem_nonZeroDivisors} it carries the nonzerodivisor \(X - b\) to a
  nonzerodivisor, which is \(u \otimes 1 - 1 \otimes b\).
\end{proof}

\section{Vanishing propagation via residue fields}
\label{pa-sec:nakayama_vanishing}

Let \(R\) be a commutative ring. For a prime ideal \(\mathfrak p \subseteq R\) write
\(R_{\mathfrak p}\) for the localisation, \(\kappa(\mathfrak p) =
R_{\mathfrak p}/\mathfrak p R_{\mathfrak p}\) for the residue field, and
\(M_{\mathfrak p}\) for the localisation of an \(R\)-module \(M\). The \emph{fibre} of
\(M\) at \(\mathfrak p\) is the \(\kappa(\mathfrak p)\)-vector space
\(M \otimes_R \kappa(\mathfrak p)\), and the \emph{support} of \(M\) is
\(\operatorname{Supp} M = \{\mathfrak p \text{ prime} : M_{\mathfrak p} \neq 0\}\).
This section proves that for a finitely generated module the vanishing of all fibres
forces the module to vanish, and that the fibre-vanishing locus is open in
\(\Spec R\). Applied with \(M = H^1\) of the relative curve
(Chapter~\ref{pa-chap:Cohomology}), these are the vanishing-propagation and openness steps
of the exchange-property analysis for the Picard functor.

\begin{lemma}[A module with vanishing localisations vanishes]
  \label{pa-lem:zero_of_forall_localization_zero}
  \lean{AlgebraicJacobian.RigidEngine.subsingleton_of_forall_subsingleton_localizedModule}
  \leanok
  \dcref{custom}

  Let \(M\) be an \(R\)-module with \(M_{\mathfrak p} = 0\) for every prime
  \(\mathfrak p\). Then \(M = 0\). (No finiteness hypothesis is needed.)
\end{lemma}
\begin{proof}
  \leanok
  Suppose \(m \in M\) is nonzero. Its annihilator
  \(\operatorname{Ann}(m) = \{r \in R : rm = 0\}\) is an ideal not containing \(1\),
  hence contained in a maximal ideal \(\mathfrak p\). If \(m/1 = 0\) in
  \(M_{\mathfrak p}\), some \(s \notin \mathfrak p\) satisfies \(sm = 0\),
  contradicting \(\operatorname{Ann}(m) \subseteq \mathfrak p\); so
  \(M_{\mathfrak p} \neq 0\).
\end{proof}

\begin{lemma}[The Nakayama fibre criterion]
  \label{pa-lem:fibre_vanishing_iff_localization_zero}
  \lean{AlgebraicJacobian.RigidEngine.subsingleton_residueField_tensor_iff_notMem_support}
  \leanok
  \dcref{custom}

  Let \(M\) be a finitely generated \(R\)-module and \(\mathfrak p\) a prime. Then
  \(M \otimes_R \kappa(\mathfrak p) = 0\) if and only if \(M_{\mathfrak p} = 0\), i.e.\
  \(\mathfrak p \notin \operatorname{Supp} M\).
\end{lemma}
\begin{proof}
  \leanok
  Localisation commutes with the tensor product and \(\kappa(\mathfrak p)\) is an
  \(R_{\mathfrak p}\)-module, so
  \(M \otimes_R \kappa(\mathfrak p) \cong M_{\mathfrak p} \otimes_{R_{\mathfrak p}}
  \kappa(\mathfrak p) \cong M_{\mathfrak p}/\mathfrak p M_{\mathfrak p}\)
  (the last identification because \(\kappa(\mathfrak p) =
  R_{\mathfrak p}/\mathfrak p R_{\mathfrak p}\) and \(N \otimes_A A/I \cong N/IN\)).
  If \(M_{\mathfrak p} = 0\) the fibre certainly vanishes. Conversely suppose
  \(M_{\mathfrak p} = \mathfrak p M_{\mathfrak p}\). The module \(M_{\mathfrak p}\) is
  finitely generated over the local ring \(R_{\mathfrak p}\) (the images of generators
  of \(M\) generate), whose maximal ideal is \(\mathfrak p R_{\mathfrak p}\). If
  \(M_{\mathfrak p} \neq 0\), choose generators \(x_1, \dots, x_n\) of
  \(M_{\mathfrak p}\) with \(n \geq 1\) minimal. Expanding
  \(x_n \in \mathfrak p M_{\mathfrak p}\) in the generators gives
  \(x_n = \sum_{i \leq n} a_i x_i\) with all
  \(a_i \in \mathfrak p R_{\mathfrak p}\), whence
  \((1 - a_n) x_n = \sum_{i < n} a_i x_i\). But \(1 - a_n\) lies outside the maximal
  ideal, so it is a unit, and \(x_n\) lies in the span of \(x_1, \dots, x_{n-1}\) ---
  contradicting minimality. Hence \(M_{\mathfrak p} = 0\).
\end{proof}

\begin{theorem}[Vanishing propagation]
  \label{pa-thm:fibrewise_vanishing_propagates}
  \lean{AlgebraicJacobian.RigidEngine.subsingleton_of_forall_subsingleton_residueField_tensor}
  \leanok
  \source{kleiman-picard}

  \uses{pa-lem:fibre_vanishing_iff_localization_zero, pa-lem:zero_of_forall_localization_zero}
  \dcref{custom}
  Let \(M\) be a finitely generated \(R\)-module with
  \(M \otimes_R \kappa(\mathfrak p) = 0\) for every prime \(\mathfrak p\). Then
  \(M = 0\): fibrewise vanishing over the whole base propagates to global vanishing.
\end{theorem}
\begin{proof}
  \leanok
  By \ref{pa-lem:fibre_vanishing_iff_localization_zero} every localisation
  \(M_{\mathfrak p}\) vanishes; by \ref{pa-lem:zero_of_forall_localization_zero},
  \(M = 0\).
\end{proof}

\begin{lemma}[The fibre-vanishing locus is the complement of the support]
  \label{pa-lem:fibre_vanishing_locus_compl_support}
  \lean{AlgebraicJacobian.RigidEngine.setOf_subsingleton_residueField_tensor_eq_compl_support}
  \leanok
  \dcref{custom}

  \uses{pa-lem:fibre_vanishing_iff_localization_zero}
  For a finitely generated \(R\)-module \(M\),
  \[
    \{\mathfrak p \in \Spec R : M \otimes_R \kappa(\mathfrak p) = 0\}
    \;=\; \Spec R \smallsetminus \operatorname{Supp} M .
  \]
\end{lemma}
\begin{proof}
  \leanok
  A pointwise restatement of \ref{pa-lem:fibre_vanishing_iff_localization_zero}.
\end{proof}

\begin{theorem}[Openness of the fibre-vanishing locus]
  \label{pa-thm:isOpen_fibre_vanishing}
  \lean{AlgebraicJacobian.RigidEngine.isOpen_setOf_subsingleton_residueField_tensor}
  \leanok
  \source{kleiman-picard}

  \uses{pa-lem:fibre_vanishing_locus_compl_support}
  \dcref{custom}
  For a finitely generated \(R\)-module \(M\), the fibre-vanishing locus
  \(\{\mathfrak p : M \otimes_R \kappa(\mathfrak p) = 0\}\) is open in \(\Spec R\).
\end{theorem}
\begin{proof}
  \leanok
  By \ref{pa-lem:fibre_vanishing_locus_compl_support} the locus is the complement of
  \(\operatorname{Supp} M\), so it suffices to prove
  \(\operatorname{Supp} M = V(\operatorname{Ann} M)\) for \(M\) generated by
  \(x_1, \dots, x_n\), a closed set. If
  \(\operatorname{Ann} M \not\subseteq \mathfrak p\), pick
  \(s \in \operatorname{Ann} M \smallsetminus \mathfrak p\): then \(s\) kills \(M\), so
  \(M_{\mathfrak p} = 0\). Conversely if \(M_{\mathfrak p} = 0\), then each
  \(x_i/1 = 0\) in \(M_{\mathfrak p}\), so \(s_i x_i = 0\) for some
  \(s_i \notin \mathfrak p\); the product \(s = s_1 \cdots s_n\) lies outside the prime
  \(\mathfrak p\) and kills every generator, so
  \(s \in \operatorname{Ann} M \smallsetminus \mathfrak p\). Thus a prime lies off the
  support exactly when it fails to contain \(\operatorname{Ann} M\), i.e.\
  \(\operatorname{Supp} M = V(\operatorname{Ann} M)\). (The second computation is the
  classical localisation step: a finitely generated module vanishing at \(\mathfrak p\)
  vanishes on the basic open neighbourhood \(D(s)\).)
\end{proof}

\section{Split rigidity of a surjective two-term complex}
\label{pa-sec:split_rigidity}

Throughout this section \(R\) is a commutative ring and \(\delta : A \to B\) a map of
\(R\)-modules --- a two-term complex, with
\(H^0 := \ker \delta\) and \(H^1 := B/\delta(A)\) its cohomology. In the intended
application (Chapter~\ref{pa-chap:Cohomology}) \(\delta\) is the restriction-difference map
of an affine two-cover and \(H^1\) is degree-one cohomology; everything here is pure
module algebra. The theme is \emph{rigidity}: once \(H^1 = 0\) --- that is, once
\(\delta\) is surjective --- the formation of \(H^0\) behaves as if the complex were
split. Without the surjectivity, kernels do not commute with tensor products: for
\(\delta = p \cdot (-) : \Z \to \Z\) with \(p\) prime, \(\ker \delta = 0\), yet after
\(\otimes\, \Z/p\) the map becomes zero and acquires the kernel \(\Z/p\).

\begin{lemma}[Right-exactness of the tensor product]
  \label{pa-lem:tensor_right_exact}
  \lean{rTensor_exact, lTensor_exact, LinearMap.rTensor_surjective,
    LinearMap.lTensor_surjective}
  \mathlibok
  \dcref{custom}

  Let \(M' \to M \to M'' \to 0\) be an exact sequence of \(R\)-modules and \(P\) any
  \(R\)-module. Then the sequences
  \(M' \otimes P \to M \otimes P \to M'' \otimes P \to 0\) and
  \(P \otimes M' \to P \otimes M \to P \otimes M'' \to 0\) are exact.
\end{lemma}

\begin{lemma}[A flat cokernel keeps the kernel inclusion injective under tensoring]
  \label{pa-lem:lTensor_injective_of_flat}
  \lean{LinearMap.lTensor_injective_of_exact_of_flat}
  \mathlibok
  \dcref{custom}

  Let \(0 \to K \xrightarrow{\iota} A \xrightarrow{\delta} B \to 0\) be a short exact
  sequence of \(R\)-modules with \(B\) flat. Then for every \(R\)-module \(P\) the map
  \(\mathrm{id}_P \otimes \iota : P \otimes K \to P \otimes A\) is injective.
\end{lemma}

\begin{lemma}[A vanishing cokernel means surjectivity]
  \label{pa-lem:rigid_surjective_of_coker_zero}
  \lean{AlgebraicJacobian.RigidEngine.surjective_of_subsingleton_quotient}
  \leanok
  \dcref{custom}

  If \(H^1 = B/\delta(A) = 0\), then \(\delta\) is surjective.
\end{lemma}
\begin{proof}
  \leanok
  Every \(b \in B\) has class zero in the zero quotient, i.e.\ lies in \(\delta(A)\).
\end{proof}

\begin{lemma}[The cokernel of a surjective complex base-changes to zero]
  \label{pa-lem:rigid_coker_tensor_zero}
  \lean{AlgebraicJacobian.RigidEngine.rTensor_surjective,
    AlgebraicJacobian.RigidEngine.subsingleton_coker_rTensor}
  \leanok
  \dcref{custom}

  \uses{pa-lem:rigid_surjective_of_coker_zero}
  Let \(\delta : A \to B\) be surjective (equivalently, \(H^1 = 0\);
  \ref{pa-lem:rigid_surjective_of_coker_zero}). Then for every \(R\)-module \(P\) the map
  \(\delta \otimes \mathrm{id}_P : A \otimes P \to B \otimes P\) is surjective;
  equivalently, the cokernel of the base-changed complex vanishes for every coefficient
  module: \(H^1\) base-changes to zero.
\end{lemma}
\begin{proof}
  \leanok
  \(B \otimes P\) is generated by the pure tensors \(b \otimes p\); writing
  \(b = \delta(a)\) exhibits each as
  \((\delta \otimes \mathrm{id})(a \otimes p)\). The image, a submodule containing a
  generating set, is everything.
\end{proof}

\begin{theorem}[Base change of \(H^0\) with module coefficients]
  \label{pa-thm:rigid_ker_tensor}
  \lean{AlgebraicJacobian.RigidEngine.kerRTensorToKer,
    AlgebraicJacobian.RigidEngine.kerRTensorEquiv}
  \leanok
  \source{kleiman-picard}

  \uses{pa-lem:rigid_surjective_of_coker_zero, pa-lem:tensor_right_exact,
    pa-lem:lTensor_injective_of_flat}
  \dcref{custom}
  Let \(\delta : A \to B\) be surjective (equivalently \(H^1 = 0\)) with \(B\) flat, and
  set \(K = \ker \delta\). For every \(R\)-module \(P\), the natural map
  \(K \otimes P \to \ker(\delta \otimes \mathrm{id}_P)\) induced by the inclusion
  \(K \subseteq A\) is an isomorphism: the formation of \(H^0\) commutes with
  \(\otimes\, P\) for \emph{every} \(P\).
\end{theorem}
\begin{proof}
  \leanok
  Since \(\delta\) is surjective, \(0 \to K \xrightarrow{\iota} A \xrightarrow{\delta}
  B \to 0\) is exact. The natural map is well defined:
  \((\delta \otimes \mathrm{id}) \circ (\iota \otimes \mathrm{id}) =
  (\delta \iota) \otimes \mathrm{id} = 0\), so \(\iota \otimes \mathrm{id}\) lands in
  the kernel. \emph{Injectivity}: the cokernel \(B\) of \(\iota\) is flat, so
  \(\iota \otimes \mathrm{id}_P\) remains injective
  (\ref{pa-lem:lTensor_injective_of_flat}), and the natural map is its corestriction.
  \emph{Surjectivity onto the kernel}: by right-exactness
  (\ref{pa-lem:tensor_right_exact}) applied to the exact \(K \to A \to B \to 0\), the
  sequence \(K \otimes P \to A \otimes P \to B \otimes P\) is exact at \(A \otimes P\);
  an element of \(\ker(\delta \otimes \mathrm{id})\) therefore comes from
  \(K \otimes P\). (Flatness is not used in this half.)
\end{proof}

\begin{theorem}[Base change of \(H^0\) along a ring map]
  \label{pa-thm:rigid_ker_baseChange}
  \lean{AlgebraicJacobian.RigidEngine.kerBaseChangeToKer,
    AlgebraicJacobian.RigidEngine.kerBaseChangeEquiv}
  \leanok
  \source{kleiman-picard}

  \uses{pa-thm:rigid_ker_tensor}
  \dcref{custom}
  Let \(\delta : A \to B\) be surjective with \(B\) flat, and let \(R'\) be a
  commutative \(R\)-algebra. Write \(\delta_{R'} : R' \otimes_R A \to R' \otimes_R B\)
  for the base-changed \(R'\)-linear map. Then the natural \(R'\)-linear map
  \[
    R' \otimes_R \ker \delta \;\longrightarrow\; \ker \delta_{R'}
  \]
  is an isomorphism: the formation of \(H^0\) commutes with every base change
  \(R \to R'\).
\end{theorem}
\begin{proof}
  \leanok
  The underlying additive map is that of \ref{pa-thm:rigid_ker_tensor} with \(P = R'\)
  (with the tensor factors flipped), hence bijective; and it is \(R'\)-linear because
  it is induced by \(\mathrm{id}_{R'} \otimes \iota\), which commutes with the
  \(R'\)-action carried by the left tensor factor.
\end{proof}

\begin{theorem}[Flatness of \(H^0\): the splitting argument]
  \label{pa-thm:rigid_flat_ker}
  \lean{AlgebraicJacobian.RigidEngine.flat_ker}
  \leanok
  \dcref{custom}

  Let \(\delta : A \to B\) be surjective with \(A\) flat and \(B\) projective. Then
  \(K = \ker \delta\) is flat.
\end{theorem}
\begin{proof}
  \leanok
  Projectivity of \(B\) lifts \(\mathrm{id}_B\) through the surjection \(\delta\):
  there is \(\sigma : B \to A\) with \(\delta \sigma = \mathrm{id}_B\), so the complex
  splits. The map \(\pi := \mathrm{id}_A - \sigma\delta\) lands in \(K\)
  (\(\delta\pi = \delta - \delta\sigma\delta = 0\)) and restricts to the identity on
  \(K\) (if \(\delta x = 0\) then \(\pi x = x\)); moreover
  \(A = K \oplus \sigma(B)\): every \(a\) is \(\pi a + \sigma(\delta a)\), and if
  \(x = \sigma b \in K \cap \sigma(B)\) then \(0 = \delta x = b\), so \(x = 0\). A
  direct summand of a flat module is flat: for an injective map \(M' \to M\), the map
  \(K \otimes M' \to K \otimes M\) is a direct summand of the injective map
  \(A \otimes M' \to A \otimes M\) (the tensor product distributes over direct sums),
  hence injective. So \(K\) is flat.
\end{proof}

\begin{lemma}[Finite modules over a Noetherian ring are finitely presented]
  \label{pa-lem:finitePresentation_of_finite}
  \lean{Module.finitePresentation_of_finite}
  \mathlibok
  \dcref{custom}

  Over a Noetherian ring \(R\), every finitely generated module is finitely presented:
  the kernel of any surjection \(R^n \to M\) is again finitely generated.
\end{lemma}

\begin{lemma}[Finitely presented flat modules are projective]
  \label{pa-lem:projective_of_finitePresentation}
  \lean{Module.Flat.projective_of_finitePresentation}
  \mathlibok
  \dcref{custom}

  A finitely presented flat module over a commutative ring is projective.
\end{lemma}

\begin{theorem}[\(H^0\) is finite projective]
  \label{pa-thm:rigid_projective_ker}
  \lean{AlgebraicJacobian.RigidEngine.projective_ker}
  \leanok
  \source{kleiman-picard}

  \uses{pa-thm:rigid_flat_ker, pa-lem:finitePresentation_of_finite,
    pa-lem:projective_of_finitePresentation}
  \dcref{custom}
  Let \(R\) be Noetherian and \(\delta : A \to B\) surjective with \(A\) flat, \(B\)
  projective, and \(K = \ker \delta\) finitely generated. Then \(K\) is a finite
  projective \(R\)-module.
\end{theorem}
\begin{proof}
  \leanok
  \(K\) is flat (\ref{pa-thm:rigid_flat_ker}) and finitely presented
  (\ref{pa-lem:finitePresentation_of_finite}); finitely presented flat modules are
  projective (\ref{pa-lem:projective_of_finitePresentation}).
\end{proof}

\begin{theorem}[Base change commutes with cokernels]
  \label{pa-thm:quot_range_baseChange}
  \lean{LinearMap.quotRangeBaseChangeEquiv, LinearMap.quotRangeBaseChangeEquiv_tmul_mk}
  \leanok
  \dcref{custom}

  \uses{pa-lem:tensor_right_exact}
  Let \(f : M \to N\) be an \(R\)-linear map and \(S\) a commutative \(R\)-algebra.
  There is an \(S\)-linear equivalence
  \[
    S \otimes_R \bigl(N/\mathrm{range}\, f\bigr) \;\cong\;
      (S \otimes_R N)\big/\mathrm{range}(f_S), \qquad
    s \otimes [n] \longmapsto [s \otimes n],
  \]
  where \(f_S = \mathrm{id}_S \otimes f\) is the base-changed map.
\end{theorem}
\begin{proof}
  \leanok
  Apply \(\mathrm{id}_S \otimes -\) to the exact sequence
  \(M \xrightarrow{f} N \xrightarrow{q} N/\mathrm{range}\, f \to 0\) (\(q\) the quotient
  map): by right-exactness (\ref{pa-lem:tensor_right_exact}) the sequence
  \(S \otimes M \to S \otimes N \to S \otimes (N/\mathrm{range}\, f) \to 0\) is exact.
  So \(\mathrm{id}_S \otimes q\) is surjective with kernel \(\mathrm{range}(f_S)\), and
  the first isomorphism theorem gives an \(S\)-linear equivalence
  \((S \otimes N)/\mathrm{range}(f_S) \cong S \otimes (N/\mathrm{range}\, f)\), inverse
  to the displayed one; on classes of pure tensors it sends \([s \otimes n]\) to
  \(s \otimes [n]\).
\end{proof}

\section{Duality for finite projective modules}
\label{pa-sec:finite_projective_duality}

For an \(R\)-module \(M\) write \(M^{\vee} = \Hom_R(M, R)\) for its dual. The
\emph{contraction map}
\[
  c_M : M^{\vee} \otimes_R N \longrightarrow \Hom_R(M, N), \qquad
  \psi \otimes n \longmapsto \bigl(x \mapsto \psi(x)\, n\bigr)
\]
is natural in both variables. This section upgrades the classical fact that \(c_M\) is
an isomorphism for finite free \(M\) to finite \emph{projective} \(M\), by a retract
argument, and derives the duality \(M \otimes P \cong \Hom(M^{\vee}, P)\) --- the form
in which the representing module of the exchange isomorphism is consumed
(Chapter~\ref{pa-chap:Cohomology}).

\begin{lemma}[Contraction is an equivalence for finite free modules]
  \label{pa-lem:dualTensorHomEquiv_free}
  \lean{dualTensorHomEquiv}
  \mathlibok
  \dcref{custom}

  For a finite free \(R\)-module \(F\) and every \(R\)-module \(N\), the contraction
  map \(c_F : F^{\vee} \otimes N \to \Hom(F, N)\) is an isomorphism.
\end{lemma}

\begin{theorem}[Contraction is an equivalence for finite projective modules]
  \label{pa-thm:dualTensorHom_projective}
  \lean{AlgebraicJacobian.RigidEngine.bijective_dualTensorHom_of_projective,
    AlgebraicJacobian.RigidEngine.dualTensorHomEquivProjective}
  \leanok
  \dcref{custom}

  \uses{pa-lem:dualTensorHomEquiv_free}
  For a finite projective \(R\)-module \(M\) and every \(R\)-module \(N\), the
  contraction map \(c_M : M^{\vee} \otimes N \to \Hom(M, N)\) is an isomorphism.
\end{theorem}
\begin{proof}
  \leanok
  A finite projective module is a retract of a finite free one: generators give a
  surjection \(f : F := R^n \to M\), and projectivity lifts \(\mathrm{id}_M\) through
  \(f\) to \(g : M \to F\) with \(f \circ g = \mathrm{id}_M\). Contraction is natural in
  the module variable; concretely, on pure tensors one checks the two identities
  \[
    c_M \circ (g^{\vee} \otimes \mathrm{id}_N) = (- \circ g) \circ c_F,
    \qquad
    (- \circ f) \circ c_M = c_F \circ (f^{\vee} \otimes \mathrm{id}_N),
  \]
  both sides of the first sending \(\psi \otimes n\) to \(x \mapsto \psi(g x) n\), and
  of the second to \(x \mapsto \psi(f x) n\). Define
  \(\Phi : \Hom(M, N) \to M^{\vee} \otimes N\) by
  \(\Phi(h) := (g^{\vee} \otimes \mathrm{id}_N)\bigl(c_F^{-1}(h \circ f)\bigr)\), using
  \ref{pa-lem:dualTensorHomEquiv_free}. Then
  \(c_M(\Phi h) = c_F\bigl(c_F^{-1}(h \circ f)\bigr) \circ g = h \circ f \circ g = h\)
  by the first identity; and
  \(\Phi(c_M y) = (g^{\vee} \otimes \mathrm{id})\bigl((f^{\vee} \otimes \mathrm{id})
  y\bigr) = \bigl((g^{\vee} \circ f^{\vee}) \otimes \mathrm{id}\bigr) y = y\) by the
  second identity and \(g^{\vee} \circ f^{\vee} = (f \circ g)^{\vee} =
  \mathrm{id}\). So \(\Phi\) is a two-sided inverse of \(c_M\).
\end{proof}

\begin{theorem}[Duality: \(M \otimes P \cong \Hom(M^{\vee}, P)\)]
  \label{pa-thm:tensorDualHom_projective}
  \lean{AlgebraicJacobian.RigidEngine.tensorDualHomEquiv}
  \leanok
  \source{kleiman-picard}

  \uses{pa-thm:dualTensorHom_projective}
  \dcref{custom}
  For a finite projective \(R\)-module \(M\) and every \(R\)-module \(P\), there is a
  natural equivalence
  \[
    M \otimes_R P \;\cong\; \Hom_R(M^{\vee}, P) :
  \]
  the functor \(P \mapsto M \otimes P\) is represented, in dual-Hom form, by
  \(M^{\vee}\).
\end{theorem}
\begin{proof}
  \leanok
  Keep \(f : F \to M\), \(g : M \to F\) with \(f g = \mathrm{id}_M\) as in the proof of
  \ref{pa-thm:dualTensorHom_projective}. First, \(M^{\vee}\) is again finite projective:
  \(f^{\vee} : M^{\vee} \to F^{\vee}\) and \(g^{\vee} : F^{\vee} \to M^{\vee}\) satisfy
  \(g^{\vee} \circ f^{\vee} = \mathrm{id}\), so \(M^{\vee}\) is a retract --- hence a
  direct summand and a quotient --- of \(F^{\vee} \cong R^n\). Second, \(M\) is
  reflexive: the evaluation map \(\mathrm{ev}_M : M \to M^{\vee\vee}\),
  \(\mathrm{ev}(x)(\psi) = \psi(x)\), is an isomorphism. For finite free \(F\) this is
  the dual-basis computation; in general, naturality
  \(\mathrm{ev}_F \circ u = u^{\vee\vee} \circ \mathrm{ev}_M\) (for \(u : M \to F\))
  makes \(\beta := f \circ \mathrm{ev}_F^{-1} \circ g^{\vee\vee}\) a two-sided inverse:
  \(\beta \circ \mathrm{ev}_M = f \circ \mathrm{ev}_F^{-1} \circ \mathrm{ev}_F \circ g
  = f g = \mathrm{id}\), and \(\mathrm{ev}_M \circ \beta = f^{\vee\vee} \circ
  \mathrm{ev}_F \circ \mathrm{ev}_F^{-1} \circ g^{\vee\vee} = (f g)^{\vee\vee} =
  \mathrm{id}\). Now compose:
  \[
    M \otimes P \;\cong\; M^{\vee\vee} \otimes P
      = (M^{\vee})^{\vee} \otimes P \;\cong\; \Hom(M^{\vee}, P),
  \]
  the first map by \(\mathrm{ev}_M\), the second by
  \ref{pa-thm:dualTensorHom_projective} applied to the finite projective module
  \(M^{\vee}\).
\end{proof}

\section{Transport of kernels and of chart-lattice finiteness}
\label{pa-sec:aeval_transport}

Two small transport toolkits consumed by the sheaf-level assembly of the rigid engine
(Chapter~\ref{pa-chap:Cohomology}). For an endomorphism \(e\) of an \(S\)-module \(M\),
write \(M_e\) for \(M\) viewed as an \(S[t]\)-module with \(t\) acting through \(e\) ---
the \emph{chart lattice} of \(e\); ``\(M_e\) is finite over \(S[t]\)'' is the finiteness
hypothesis under which the coherence theorems of two-lattice pairs fire
(\ref{pa-thm:twoLattice_coh1}, \ref{pa-thm:twoLattice_coh0}).

\begin{lemma}[Kernel and surjectivity transport along a commuting square]
  \label{pa-lem:kerCongr}
  \lean{AlgebraicJacobian.RigidEngine.kerCongr,
    AlgebraicJacobian.RigidEngine.kerCongr_apply_coe,
    AlgebraicJacobian.RigidEngine.surjective_congr}
  \leanok
  \dcref{custom}

  Let \(\delta : A \to B\) and \(\delta' : A' \to B'\) be \(S\)-linear maps and
  \(e_1 : A \cong A'\), \(e_2 : B \cong B'\) \(S\)-linear equivalences with
  \(e_2 \circ \delta = \delta' \circ e_1\). Then \(e_1\) restricts to an \(S\)-linear
  equivalence \(\ker \delta \cong \ker \delta'\), and \(\delta\) is surjective if and
  only if \(\delta'\) is.
\end{lemma}
\begin{proof}
  \leanok
  If \(\delta a = 0\) then \(\delta'(e_1 a) = e_2(\delta a) = 0\), so \(e_1\) maps
  \(\ker\delta\) into \(\ker\delta'\); symmetrically \(e_1^{-1}\) maps
  \(\ker\delta'\) into \(\ker\delta\), using \(\delta = e_2^{-1} \circ \delta' \circ
  e_1\). The two restrictions are mutually inverse and \(S\)-linear. For surjectivity:
  given \(b' \in B'\), pick \(a\) with \(\delta a = e_2^{-1} b'\); then
  \(\delta'(e_1 a) = e_2(\delta a) = b'\); the converse is symmetric.
\end{proof}

\begin{lemma}[Chart-lattice transport along an equivariant equivalence]
  \label{pa-lem:aevalCongr}
  \lean{AlgebraicJacobian.RigidEngine.aevalCongr,
    AlgebraicJacobian.RigidEngine.moduleFinite_aeval'_of_linearEquiv}
  \leanok
  \dcref{custom}

  Let \(e, e'\) be endomorphisms of the \(S\)-modules \(M, M'\) and
  \(\varphi : M \cong M'\) an \(S\)-linear equivalence intertwining them,
  \(\varphi \circ e = e' \circ \varphi\). Then \(\varphi\) is an isomorphism of
  \(S[t]\)-modules \(M_e \cong M'_{e'}\); in particular, if \(M_e\) is a finite
  \(S[t]\)-module then so is \(M'_{e'}\).
\end{lemma}
\begin{proof}
  \leanok
  The \(S[t]\)-action on \(M_e\) is generated over the \(S\)-action by the action of
  \(t\), which is \(e\); an \(S\)-linear map intertwining \(e\) with \(e'\) therefore
  commutes with the action of every polynomial, i.e.\ is \(S[t]\)-linear. Being
  bijective, \(\varphi\) is an \(S[t]\)-isomorphism, and finiteness transports along it
  (the image of a finite generating set generates).
\end{proof}

\begin{lemma}[Chart-lattice finiteness from a finite ring map]
  \label{pa-lem:moduleFinite_aeval_ringHom}
  \lean{AlgebraicJacobian.RigidEngine.aeval_apply_of_ringHom,
    AlgebraicJacobian.RigidEngine.moduleFinite_aeval'_of_ringHom_finite}
  \leanok
  \dcref{custom}

  Let \(k\) be a commutative ring, \(A\) a commutative ring which is also a
  \(k\)-module, and \(\varphi : k[t] \to A\) a \emph{finite} ring homomorphism. Let
  \(e\) be a \(k\)-linear endomorphism of \(A\) acting as multiplication by
  \(\varphi(t)\), and suppose the \(k\)-scalar action is multiplication by the constants,
  \(c \bullet a = \varphi(c) \cdot a\). Then the polynomial action of \(p \in k[t]\) on
  \(A_e\) is multiplication by \(\varphi(p)\), and \(A_e\) is a finite \(k[t]\)-module.
\end{lemma}
\begin{proof}
  \leanok
  For the action formula, induct on the monomials of \(p\): the action of \(c\,t^n\) on
  \(a\) is \(c \bullet e^n(a) = \varphi(c)\,\varphi(t)^n\, a = \varphi(c\,t^n)\, a\),
  using that \(e^n\) is multiplication by \(\varphi(t)^n\); sums of monomials act by
  sums. Consequently the \(k[t]\)-module \(A_e\) is exactly \(A\) viewed as a
  \(k[t]\)-module through \(\varphi\), which is finite by hypothesis: a finite
  generating set of \(A\) over \(k[t]\) through \(\varphi\) generates \(A_e\).
\end{proof}

\begin{theorem}[Base change of chart-lattice finiteness]
  \label{pa-thm:moduleFinite_aeval_baseChange}
  \lean{AlgebraicJacobian.RigidEngine.aeval_baseChange_map_apply_one_tmul,
    AlgebraicJacobian.RigidEngine.moduleFinite_aeval'_baseChange}
  \leanok
  \dcref{custom}

  Let \(k\) be a commutative ring, \(R\) a commutative \(k\)-algebra, \(M\) a
  \(k\)-module and \(e\) a \(k\)-linear endomorphism of \(M\). If \(M_e\) is a finite
  \(k[t]\)-module, then \((R \otimes_k M)_{\mathrm{id} \otimes e}\) is a finite
  \(R[t]\)-module.
\end{theorem}
\begin{proof}
  \leanok
  First, base change of the polynomial action on pure tensors: for \(p \in k[t]\) and
  \(m \in M\),
  \[
    \bigl(p^R\bigr) \cdot \bigl(1 \otimes m\bigr) \;=\; 1 \otimes (p \cdot m),
  \]
  where \(p^R \in R[t]\) is \(p\) with coefficients pushed into \(R\) --- by induction on
  monomials, since \((\mathrm{id} \otimes e)^n (1 \otimes m) = 1 \otimes e^n(m)\) and
  the \(k\)-scalars pass through the tensor sign. Now let \(z_1, \dots, z_r\) generate
  \(M_e\) over \(k[t]\); we claim the \(1 \otimes z_i\) generate
  \((R \otimes_k M)_{\mathrm{id}\otimes e}\) over \(R[t]\). Every element of
  \(R \otimes_k M\) is a finite sum of tensors \(r \otimes m\), and
  \(r \otimes m = (r\,t^0) \cdot (1 \otimes m)\) with \(r\,t^0 \in R[t]\) acting as the
  constant \(r\); so it suffices to reach the \(1 \otimes m\). Write
  \(m = \sum_i p_i \cdot z_i\) with \(p_i \in k[t]\); then by the displayed formula and
  additivity, \(1 \otimes m = \sum_i p_i^R \cdot (1 \otimes z_i)\), an \(R[t]\)-linear
  combination of the claimed generators.
\end{proof}

\section{The point-fiber ideal of an \'etale-coordinatised algebra}
\label{pa-sec:point_fiber_ideal}

The diagonal-ideal package of Section~\ref{pa-sec:diagonal_ideal} computes the kernel of the
multiplication \(B \otimes_{k[t]} B \to B\) after localising at the idempotent lift. Here
the second tensor factor is specialised along a point: \(F\) a second
\(k[t]\)-coordinatised algebra receiving \(B\) through a \(k[t]\)-algebra homomorphism
\(c \colon B \to F\) (in the application, \(F\) is the section ring of an affine field test
and \(c\) the pullback along a point of the curve). Everything mirrors the diagonal
computation, with the idempotent input \emph{pushed} from the frozen diagonal package — no
unramifiedness of \(F\) is invoked.

Throughout, \(k\) is a commutative ring, \(B\) and \(F\) are \(k\)-algebras carrying
compatible \(k[t]\)-algebra structures, \(u_B := t \cdot 1_B\) and \(u_F := t \cdot 1_F\)
denote the coordinates, and \(c \colon B \to F\) is a \(k[t]\)-algebra homomorphism (so
\(c(u_B) = u_F\)).

\begin{definition}[The point generator and the point evaluations]
  \label{pa-def:pointGen_pointEv}
  \lean{AlgebraicJacobian.Diagonal.pointGen, AlgebraicJacobian.Diagonal.map_coord,
    AlgebraicJacobian.Diagonal.pointEv, AlgebraicJacobian.Diagonal.pointEv_tmul,
    AlgebraicJacobian.Diagonal.pointEv_pointGen, AlgebraicJacobian.Diagonal.pointEvTwo,
    AlgebraicJacobian.Diagonal.pointBaseChange,
    AlgebraicJacobian.Diagonal.pointBaseChange_surjective,
    AlgebraicJacobian.Diagonal.pointEvTwo_pointBaseChange,
    AlgebraicJacobian.Diagonal.mapRight, AlgebraicJacobian.Diagonal.mapRight_diagGen,
    AlgebraicJacobian.Diagonal.mapRightTwo,
    AlgebraicJacobian.Diagonal.pointBaseChange_mapRight}
  \leanok
  \dcref{custom}

  \uses{pa-def:diagonalSetup}
  The \emph{point generator} is \(g := u_B \otimes 1 - 1 \otimes u_F \in B \otimes_k F\);
  the \emph{point evaluation} is the \(k\)-algebra map
  \(\varepsilon \colon B \otimes_k F \to F\), \(x \otimes \lambda \mapsto c(x)\lambda\),
  which kills \(g\). Auxiliary to them: the base-change surjection
  \(\beta \colon B \otimes_k F \to B \otimes_{k[t]} F\); the \(k[t]\)-level evaluation
  \(\varepsilon_2 \colon B \otimes_{k[t]} F \to F\) with
  \(\varepsilon_2 \circ \beta = \varepsilon\); and the second-factor pushes
  \(1 \otimes c \colon B \otimes_k B \to B \otimes_k F\) and
  \(B \otimes_{k[t]} B \to B \otimes_{k[t]} F\), which carry the diagonal generator to the
  point generator and intertwine the two base changes.
\end{definition}
\begin{proof}
  \leanok
  \(c(u_B) = u_F\) since \(c\) is a \(k[t]\)-algebra map, so
  \(\varepsilon(g) = c(u_B)\cdot 1 - 1 \cdot u_F = 0\), and the push of the diagonal
  generator \(u_B \otimes 1 - 1 \otimes u_B\) (Definition~\ref{pa-def:diagonalSetup}) is \(g\).
  The base change is surjective because it is the identity on pure tensors; the remaining
  identities are checked on pure tensors.
\end{proof}

\begin{theorem}[The kernel of the point base change]
  \label{pa-thm:ker_pointBaseChange}
  \lean{AlgebraicJacobian.Diagonal.coord_tmul_one_point,
    AlgebraicJacobian.Diagonal.pointBaseChange_pointGen,
    AlgebraicJacobian.Diagonal.aeval_sub_dvd_point,
    AlgebraicJacobian.Diagonal.algebraMap_tmul_sub_mem_span_point,
    AlgebraicJacobian.Diagonal.pointQRight, AlgebraicJacobian.Diagonal.pointSectionMap,
    AlgebraicJacobian.Diagonal.pointSectionMap_pointBaseChange,
    AlgebraicJacobian.Diagonal.ker_pointBaseChange}
  \leanok
  \dcref{custom}

  \uses{pa-def:pointGen_pointEv, pa-lem:diagonal_telescoping, pa-lem:diagonal_section}
  The kernel of the base-change surjection
  \(\beta \colon B \otimes_k F \to B \otimes_{k[t]} F\) is exactly the principal ideal on
  the point generator: \(\ker \beta = (g)\).
\end{theorem}
\begin{proof}
  \leanok
  Mirror of Theorem~\ref{pa-thm:ker_baseChange}. \(\beta(g) = 0\) since in
  \(B \otimes_{k[t]} F\) the two coordinates agree across the tensor (both are
  \(t \cdot (1 \otimes 1)\)), so \((g) \subseteq \ker\beta\). Conversely the telescoping
  identity holds in the point setting: for every \(p \in k[t]\),
  \(p(u_B) \otimes 1 - 1 \otimes p(u_F) \in (g)\), because a difference of evaluations is
  divisible by the difference of the points (as in Lemma~\ref{pa-lem:diagonal_telescoping}).
  Telescoping promotes the right inclusion \(\lambda \mapsto 1 \otimes \lambda\) into the
  quotient \((B \otimes_k F)/(g)\) to a \(k[t]\)-algebra map, whence a section
  \(B \otimes_{k[t]} F \to (B \otimes_k F)/(g)\) of \(\beta\) with
  \(\mathrm{section} \circ \beta = \mathrm{quotient\ map}\) (as in
  Lemma~\ref{pa-lem:diagonal_section}). An element of \(\ker\beta\) is then killed by the
  quotient map, i.e.\ lies in \((g)\).
\end{proof}

\begin{theorem}[The kernel of the {\(k[t]\)}-level point evaluation]
  \label{pa-thm:ker_pointEvTwo}
  \lean{AlgebraicJacobian.Diagonal.pointEvTwo_mapRightTwo,
    AlgebraicJacobian.Diagonal.tmul_sub_tmul_mem_span_mapRightTwo,
    AlgebraicJacobian.Diagonal.ker_pointEvTwo}
  \leanok
  \dcref{custom}

  \uses{pa-def:pointGen_pointEv}
  Suppose the diagonal ideal is principal:
  \(\ker\bigl(\mathrm{mul} \colon B \otimes_{k[t]} B \to B\bigr) = (e)\) for some
  \(e \in B \otimes_{k[t]} B\). Then the kernel of the \(k[t]\)-level point evaluation is
  the principal ideal on the push of \(e\):
  \[
    \ker\bigl(\varepsilon_2 \colon B \otimes_{k[t]} F \to F\bigr)
      = \bigl((1 \otimes c)(e)\bigr).
  \]
\end{theorem}
\begin{proof}
  \leanok
  The push kills into the kernel: \(\varepsilon_2\bigl((1 \otimes c)(z)\bigr) =
  c(\mathrm{mul}(z))\) on pure tensors, and \(\mathrm{mul}(e) = 0\).

  Conversely, the generators \(x \otimes 1 - 1 \otimes c(x)\) (\(x \in B\)) of the point
  ideal are pushes of the diagonal-ideal elements \(x \otimes 1 - 1 \otimes x\), each a
  multiple of \(e\) by hypothesis; hence they lie in \(((1\otimes c)(e))\). Now the section
  trick: modulo \(((1\otimes c)(e))\), every \(z \in B \otimes_{k[t]} F\) is congruent to
  \(1 \otimes \varepsilon_2(z)\) — on a pure tensor \(x \otimes y\), the difference
  \(x \otimes y - 1 \otimes c(x)y = (x \otimes 1 - 1 \otimes c(x))(1 \otimes y)\) is a
  multiple of a pushed generator. So \(\varepsilon_2(z) = 0\) forces
  \(z \equiv 1 \otimes 0 = 0\), i.e.\ \(z \in ((1\otimes c)(e))\).
\end{proof}

\begin{theorem}[The localised point-fiber ideal]
  \label{pa-thm:ker_pointEv_map_localization_eq}
  \lean{AlgebraicJacobian.Diagonal.ker_pointEv_le_sup,
    AlgebraicJacobian.Diagonal.ker_pointEv_map_localization_eq}
  \leanok
  \dcref{custom}

  \uses{pa-thm:ker_pointBaseChange, pa-thm:ker_pointEvTwo, pa-def:pointGen_pointEv,
    pa-thm:exists_idempotent_lift, pa-thm:ker_map_localization_eq}
  Let \(\tilde e \in B \otimes_k B\) be a lift of an idempotent generator of the diagonal
  ideal: suppose \(\beta_B(\tilde e)\) is idempotent and
  \(\ker(\mathrm{mul}) = (\beta_B(\tilde e))\), where
  \(\beta_B \colon B \otimes_k B \to B \otimes_{k[t]} B\) is the diagonal base change (such
  a lift exists for the \'etale coordinate by Theorem~\ref{pa-thm:exists_idempotent_lift}).
  Write \(\tilde e_F := (1 \otimes c)(\tilde e) \in B \otimes_k F\) for its push. Then for
  \emph{any} localisation \(L\) of \(B \otimes_k F\) away from \(1 - \tilde e_F\),
  \[
    (\ker \varepsilon)\, L \;=\; \bigl(g\bigr) L ,
  \]
  the extension of the point-evaluation kernel is the principal ideal on the point
  generator.
\end{theorem}
\begin{proof}
  \leanok
  Mirror of Theorem~\ref{pa-thm:ker_map_localization_eq}.

  \emph{Containment into the sum.} \(\ker\varepsilon \subseteq (g) + (\tilde e_F)\): for
  \(x \in \ker\varepsilon\), the image \(\beta(x)\) lies in \(\ker\varepsilon_2\)
  (Definition~\ref{pa-def:pointGen_pointEv}), which is principal on the pushed generator
  (Theorem~\ref{pa-thm:ker_pointEvTwo}, whose hypothesis holds with
  \(e := \beta_B(\tilde e)\), the pushes intertwining by
  Definition~\ref{pa-def:pointGen_pointEv}); lift the factorisation through the surjection
  \(\beta\) and correct by \(\ker\beta = (g)\) (Theorem~\ref{pa-thm:ker_pointBaseChange}).

  \emph{The pushed lift dies in the localisation.} \(\beta(\tilde e_F)\) is idempotent (the
  push of an idempotent), so \(\tilde e_F (1 - \tilde e_F) \in \ker\beta = (g)\); in \(L\)
  the factor \(1 - \tilde e_F\) is invertible, hence \(\tilde e_F \in (g)L\).

  Combining, \((\ker\varepsilon)L \subseteq (g)L\); the reverse containment is
  \(\varepsilon(g) = 0\) (Definition~\ref{pa-def:pointGen_pointEv}).
\end{proof}

\section{Depth and the Auslander--Buchsbaum formula}
\label{pa-sec:auslander_buchsbaum}

This section develops the depth theory of finite modules over a Noetherian local ring
--- the numerical depth, its characterisation by the vanishing of
\(\Ext^{\bullet}(\kappa, -)\), and its behaviour in short exact sequences --- and proves
the Auslander--Buchsbaum formula
\(\mathrm{pd}_R(M) + \mathrm{depth}(M) = \mathrm{depth}(R)\).

\begin{definition}[Depth of a module]
  \label{pa-def:depth}
  \lean{RingTheory.Module.depth}
  \leanok
  \source{stacks-project}

  \dcref{custom}
  Let \(R\) be a commutative ring, \(I \subseteq R\) an ideal and \(M\) an \(R\)-module.
  An \emph{\(M\)-regular sequence} is a finite sequence \(x_1, \dots, x_n \in R\) such
  that each \(x_i\) is a non-zero-divisor on \(M/(x_1, \dots, x_{i-1})M\) and
  \((x_1, \dots, x_n)M \neq M\). The \emph{\(I\)-depth} of \(M\) is the element
  \(\mathrm{depth}_I(M)\) of \(\{0, 1, 2, \dots, \infty\}\) defined by:
  \begin{enumerate}
    \item if \(IM = M\), then \(\mathrm{depth}_I(M) = \infty\);
    \item otherwise, \(\mathrm{depth}_I(M)\) is the supremum of the lengths of
      \(M\)-regular sequences contained in \(I\) (the supremum of the empty set
      being \(0\)).
  \end{enumerate}
  When \((R, \mathfrak m)\) is local, \(\mathrm{depth}(M) := \mathrm{depth}_{\mathfrak m}(M)\)
  is called simply the \emph{depth} of \(M\).
\end{definition}

\begin{definition}[Projective dimension]
  \label{pa-def:projectiveDimension}
  \lean{Module.projectiveDimension}
  \leanok
  \dcref{custom}

  For a ring \(R\) and an \(R\)-module \(M\), the \emph{projective dimension}
  \(\mathrm{pd}_R(M) \in \{-\infty\} \cup \{0, 1, 2, \dots, \infty\}\) is the categorical
  projective dimension of \(M\) in the abelian category of \(R\)-modules: the infimum of
  the \(n\) such that \(\Ext^i(M, -) = 0\) for all \(i > n\), equivalently the smallest
  length of a projective resolution of \(M\); it is \(-\infty\) exactly when \(M = 0\).
\end{definition}

\begin{lemma}[The covariant long exact sequence of \(\Ext\)]
  \label{pa-lem:ext_covariant_les}
  \lean{CategoryTheory.Abelian.Ext.covariant_sequence_exact₁,
    CategoryTheory.Abelian.Ext.covariant_sequence_exact₂,
    CategoryTheory.Abelian.Ext.covariant_sequence_exact₃,
    CategoryTheory.Abelian.Ext.postcomp_mk₀_injective_of_mono}
  \mathlibok
  \dcref{custom}

  \uses{pa-lem:ext_contravariant_les}
  Let \(0 \to A \xrightarrow{f} B \xrightarrow{g} C \to 0\) be a short exact sequence in
  an abelian category with \(\Ext\)-groups, with extension class
  \(\delta \in \Ext^1(C, A)\) (\ref{pa-lem:ext_contravariant_les}), and let \(Y\) be any
  object. For every \(n \geq 0\) the sequence of abelian groups
  \[
    \Ext^n(Y, A) \xrightarrow{f_*} \Ext^n(Y, B) \xrightarrow{g_*} \Ext^n(Y, C)
      \xrightarrow{(-) \cdot \delta} \Ext^{n+1}(Y, A) \xrightarrow{f_*} \Ext^{n+1}(Y, B)
  \]
  is exact at each inner term, where \(f_*, g_*\) are postcomposition (covariant
  functoriality of \(\Ext(Y, -)\)) and \((-) \cdot \delta\) is the Yoneda product with
  the extension class. Moreover, if \(f : A \to B\) is a monomorphism then
  postcomposition \(f_* : \Ext^0(Y, A) \to \Ext^0(Y, B)\) is injective (in degree zero
  \(\Ext^0 = \Hom\), and postcomposition with a mono is injective on \(\Hom\)).
\end{lemma}

\begin{lemma}[Annihilator elements kill \(\Ext\)]
  \label{pa-lem:ext_annihilator_smul}
  \lean{RingTheory.Module.ext_smul_eq_zero_of_mem_annihilator}
  \leanok
  \dcref{custom}

  Let \(R\) be a commutative ring, \(N, M\) two \(R\)-modules, \(i \geq 0\), and
  \(x \in \mathrm{Ann}_R(N)\). Then \(x\) annihilates \(\Ext^i_R(N, M)\): for every
  class \(e \in \Ext^i_R(N, M)\) we have \(x \cdot e = 0\).
\end{lemma}
\begin{proof}
  \leanok
  The \(R\)-action on \(\Ext^i_R(N, M)\) can be computed through the first argument: by
  \(R\)-bilinearity of the Yoneda composition and functoriality,
  \(x \cdot e = (x \cdot \mathrm{id}_N)^*(e)\), the precomposition of \(e\) with the
  endomorphism ``multiplication by \(x\)'' of \(N\). Since \(x \in \mathrm{Ann}_R(N)\),
  the endomorphism \(x \cdot \mathrm{id}_N\) is the zero map of \(N\), and
  precomposition with the zero morphism is the zero map on \(\Ext\). Hence
  \(x \cdot e = 0\).
\end{proof}

\begin{theorem}[Depth via \(\Ext\)]
  \label{pa-thm:depth_ext}
  \lean{RingTheory.Module.depth_eq_smallest_ext_index}
  \leanok
  \source{stacks-project}

  \uses{pa-def:depth, pa-lem:ext_annihilator_smul, pa-lem:ext_covariant_les}
  \dcref{custom}
  Let \((R, \mathfrak m)\) be a Noetherian local ring with residue field
  \(\kappa = R/\mathfrak m\) and let \(M\) be a nonzero finite \(R\)-module. For every
  \(n \in \mathbb{N}\),
  \[
    n \leq \mathrm{depth}(M)
    \iff
    \Ext^i_R(\kappa, M) = 0 \text{ for all } i < n.
  \]
  Equivalently, \(\mathrm{depth}(M)\) is the smallest \(i\) such that
  \(\Ext^i_R(\kappa, M) \neq 0\).
\end{theorem}
\begin{proof}
  \leanok
  First note two facts used repeatedly. (a) Since \(M\) is nonzero and finite,
  Nakayama's lemma rules out \(\mathfrak m M = M\), so \(\mathrm{depth}(M)\) is the
  supremum of Definition~\ref{pa-def:depth}; moreover, for the same reason, any
  \emph{prefix} of an \(M\)-regular sequence contained in \(\mathfrak m\) is again
  \(M\)-regular (weak regularity passes to prefixes, and the nondegeneracy condition
  \((x_1,\dots,x_j)M \neq M\) holds because \((x_1,\dots,x_j) \subseteq \mathfrak m\) and
  \(\mathfrak m M \neq M\)). So if \(n < \mathrm{depth}(M)\) there is an \(M\)-regular
  sequence in \(\mathfrak m\) of length exactly \(n\). (b) \(\mathrm{Ann}_R(\kappa) =
  \mathfrak m\), so every \(x \in \mathfrak m\) annihilates every \(\Ext^i(\kappa, -)\)
  by \ref{pa-lem:ext_annihilator_smul}.

  We induct on \(n\), generalising \(M\). The case \(n = 0\) is trivial (both sides
  always hold). Assume the statement for \(n\) and all nonzero finite modules.

  (\(\Rightarrow\)) Let \(n + 1 \leq \mathrm{depth}(M)\). By (a) extract an
  \(M\)-regular sequence \(x, x_1, \dots, x_n\) of length \(n+1\) in \(\mathfrak m\).
  Then \(x\) is a non-zero-divisor on \(M\), the quotient \(M/xM\) is nonzero (Nakayama,
  \(x \in \mathfrak m\)) and finite, and \(x_1, \dots, x_n\) is \((M/xM)\)-regular, so
  \(n \leq \mathrm{depth}(M/xM)\); by the inductive hypothesis
  \(\Ext^j(\kappa, M/xM) = 0\) for all \(j < n\). Now fix \(i \leq n\) and consider the
  short exact sequence \(0 \to M \xrightarrow{x} M \to M/xM \to 0\). If \(i = 0\):
  postcomposition with the monomorphism \(x : M \to M\) is injective on
  \(\Ext^0(\kappa, M)\) (\ref{pa-lem:ext_covariant_les}) and coincides with the action of
  \(x\), which is zero by (b); hence \(\Ext^0(\kappa, M) = 0\). If \(i = j + 1\) with
  \(j < n\): in the covariant long exact sequence (\ref{pa-lem:ext_covariant_les})
  \[
    \Ext^{j}(\kappa, M/xM) \xrightarrow{\ \delta\ } \Ext^{j+1}(\kappa, M)
      \xrightarrow{\ x_*\ } \Ext^{j+1}(\kappa, M),
  \]
  the map \(x_*\) (postcomposition with multiplication by \(x\)) is the action of \(x\),
  hence zero by (b); exactness makes \(\delta\) surjective, and its source vanishes, so
  \(\Ext^{j+1}(\kappa, M) = 0\).

  (\(\Leftarrow\)) Suppose \(\Ext^i(\kappa, M) = 0\) for all \(i \leq n\). From
  \(\Ext^0(\kappa, M) = \Hom_R(\kappa, M) = 0\) it follows that some
  \(x \in \mathfrak m = \mathrm{Ann}_R(\kappa)\) is a non-zero-divisor on \(M\) (for a
  finite module over a Noetherian ring, if every element of \(\mathrm{Ann}(\kappa)\)
  were a zero-divisor on \(M\), then \(\mathfrak m\) would be contained in an associated
  prime of \(M\), producing a nonzero map \(\kappa \to M\)). The quotient \(M/xM\) is
  nonzero (Nakayama) and finite. For \(j < n\), the segment
  \[
    \Ext^{j}(\kappa, M) \longrightarrow \Ext^{j}(\kappa, M/xM)
      \xrightarrow{\ \delta\ } \Ext^{j+1}(\kappa, M)
  \]
  of the covariant long exact sequence (\ref{pa-lem:ext_covariant_les}) has both outer
  terms zero (\(j < n\) and \(j + 1 \leq n\)), so \(\Ext^{j}(\kappa, M/xM) = 0\). By the
  inductive hypothesis \(n \leq \mathrm{depth}(M/xM)\); by (a) pick an
  \((M/xM)\)-regular sequence \(x_1, \dots, x_n\) of length \(n\) in \(\mathfrak m\).
  Then \(x, x_1, \dots, x_n\) is an \(M\)-regular sequence of length \(n+1\) in
  \(\mathfrak m\), whence \(n + 1 \leq \mathrm{depth}(M)\).

  Finally, the two formulations are equivalent: the displayed equivalence for all \(n\)
  says exactly that \(\mathrm{depth}(M)\) is the least index with nonvanishing
  \(\Ext(\kappa, M)\).
\end{proof}

\begin{theorem}[Depth in short exact sequences]
  \label{pa-thm:depth_ses}
  \lean{RingTheory.Module.depth_of_short_exact}
  \leanok
  \source{stacks-project}

  \uses{pa-thm:depth_ext, pa-lem:ext_covariant_les}
  \dcref{custom}
  Let \((R, \mathfrak m)\) be a Noetherian local ring and
  \(0 \to N' \to N \to N'' \to 0\) a short exact sequence of nonzero finite
  \(R\)-modules. Then, with truncated subtraction in \(\{0, 1, \dots, \infty\}\),
  \begin{enumerate}
    \item \(\mathrm{depth}(N) \geq \min\{\mathrm{depth}(N'), \mathrm{depth}(N'')\}\),
    \item \(\mathrm{depth}(N'') \geq \min\{\mathrm{depth}(N), \mathrm{depth}(N') - 1\}\),
    \item \(\mathrm{depth}(N') \geq \min\{\mathrm{depth}(N), \mathrm{depth}(N'') + 1\}\).
  \end{enumerate}
\end{theorem}
\begin{proof}
  \leanok
  Write \(\kappa = R/\mathfrak m\). In each case it suffices, by \ref{pa-thm:depth_ext}, to
  fix a natural number \(a\) below the right-hand side and prove
  \(\Ext^i(\kappa, -) = 0\) of the middle-named module for all \(i < a\); all three are
  read off the covariant long exact sequence of \(\Ext(\kappa, -)\)
  (\ref{pa-lem:ext_covariant_les}) applied to the given short exact sequence, using
  \ref{pa-thm:depth_ext} to convert the depth bounds on the other two modules into
  \(\Ext\)-vanishing.

  (1) For \(i < a \leq \min\{\mathrm{depth}(N'), \mathrm{depth}(N'')\}\): in
  \(\Ext^i(\kappa, N') \to \Ext^i(\kappa, N) \to \Ext^i(\kappa, N'')\), the outer terms
  vanish, so by exactness at the middle every class of \(\Ext^i(\kappa, N)\) lifts to
  \(\Ext^i(\kappa, N') = 0\).

  (2) For \(i < a \leq \min\{\mathrm{depth}(N), \mathrm{depth}(N') - 1\}\): in
  \(\Ext^i(\kappa, N) \to \Ext^i(\kappa, N'') \xrightarrow{\delta}
  \Ext^{i+1}(\kappa, N')\), the first term vanishes since \(i < \mathrm{depth}(N)\) and
  the last since \(i + 1 < \mathrm{depth}(N')\) (from \(i + 1 \leq a \leq
  \mathrm{depth}(N') - 1\)); exactness at the middle gives
  \(\Ext^i(\kappa, N'') = 0\).

  (3) For \(i < a \leq \min\{\mathrm{depth}(N), \mathrm{depth}(N'') + 1\}\): if
  \(i = 0\), postcomposition \(\Ext^0(\kappa, N') \to \Ext^0(\kappa, N)\) with the
  monomorphism \(N' \to N\) is injective (\ref{pa-lem:ext_covariant_les}) and lands in
  \(\Ext^0(\kappa, N) = 0\). If \(i = j + 1\), then \(j < \mathrm{depth}(N'')\) (from
  \(j + 2 \leq a \leq \mathrm{depth}(N'') + 1\)), and in
  \(\Ext^{j}(\kappa, N'') \xrightarrow{\delta} \Ext^{j+1}(\kappa, N') \to
  \Ext^{j+1}(\kappa, N)\) the outer terms vanish, so
  \(\Ext^{j+1}(\kappa, N') = 0\).
\end{proof}

\begin{lemma}[Depth is invariant under linear equivalence]
  \label{pa-lem:depth_linearEquiv}
  \lean{RingTheory.Module.depth_eq_of_linearEquiv}
  \leanok
  \dcref{custom}

  \uses{pa-def:depth}
  Let \(R\) be a commutative ring, \(I \subseteq R\) an ideal, and \(e : M \cong M'\) an
  \(R\)-linear equivalence of \(R\)-modules. Then
  \(\mathrm{depth}_I(M) = \mathrm{depth}_I(M')\).
\end{lemma}
\begin{proof}
  \leanok
  The condition \(IM = M\) transfers along \(e\) (the image of \(IM\) under \(e\) is
  \(IM'\), and \(e\) is surjective), so the two depths land in the same branch of
  Definition~\ref{pa-def:depth}. In the supremum branch, a sequence \(x_1, \dots, x_n\) is
  \(M\)-regular if and only if it is \(M'\)-regular: \(e\) intertwines the
  multiplication maps and induces equivalences on the successive quotients
  \(M/(x_1, \dots, x_j)M \cong M'/(x_1, \dots, x_j)M'\), so non-zero-divisor conditions
  and the nondegeneracy condition match term by term. Hence the two suprema agree.
\end{proof}

\begin{lemma}[Ideal action on a finite power is componentwise]
  \label{pa-lem:ideal_smul_top_pi}
  \lean{RingTheory.Module.ideal_smul_top_pi_const}
  \leanok
  \dcref{custom}

  Let \(R\) be a commutative ring, \(I \subseteq R\) an ideal, \(M\) an \(R\)-module and
  \(\iota\) a finite index type. Inside the product module \(\iota \to M\),
  \[
    I \cdot (\iota \to M) \;=\; \{\, f : \iota \to M \mid f(i) \in IM
      \text{ for all } i \,\}.
  \]
\end{lemma}
\begin{proof}
  \leanok
  \(\subseteq\): a generator \(a \cdot f\) with \(a \in I\) has every component
  \(a \cdot f(i) \in IM\), and the right-hand side is closed under sums.
  \(\supseteq\): a function \(f\) with all components in \(IM\) decomposes as the finite
  sum \(f = \sum_i \delta_i(f(i))\) of its one-hot components, where
  \(\delta_i : M \to (\iota \to M)\) places a value in slot \(i\) and zero elsewhere;
  each \(\delta_i\) is \(R\)-linear, so it carries \(IM\) into
  \(I \cdot (\iota \to M)\).
\end{proof}

\begin{lemma}[Depth of a finite power]
  \label{pa-lem:depth_pi_const}
  \lean{RingTheory.Module.depth_pi_const_eq_depth_of_nonempty}
  \leanok
  \dcref{custom}

  \uses{pa-def:depth, pa-lem:ideal_smul_top_pi, pa-lem:depth_linearEquiv}
  Let \(R\) be a commutative ring, \(I \subseteq R\) an ideal, \(M\) an \(R\)-module and
  \(\iota\) a nonempty finite index type. Then
  \(\mathrm{depth}_I(\iota \to M) = \mathrm{depth}_I(M)\). In particular
  \(\mathrm{depth}_I(R^k) = \mathrm{depth}_I(R)\) for \(k \geq 1\).
\end{lemma}
\begin{proof}
  \leanok
  The side condition agrees on both sides: by \ref{pa-lem:ideal_smul_top_pi},
  \(I \cdot (\iota \to M) = \iota \to M\) forces \(IM = M\) (evaluate a one-hot function
  at its slot) and conversely (componentwise). It remains to see that a sequence
  \(x_1, \dots, x_n \in R\) is \((\iota \to M)\)-regular iff it is \(M\)-regular; then
  the two suprema of Definition~\ref{pa-def:depth} coincide. Induct on the sequence. For
  the empty sequence, the nondegeneracy conditions \((\,)\cdot(\iota \to M) \neq
  \iota \to M\), i.e. \(\iota \to M \neq 0\), and \(M \neq 0\) agree since \(\iota\) is
  nonempty. For \(x, x_1, \dots, x_n\): multiplication by \(x\) on \(\iota \to M\) is
  componentwise, so \(x\) is a non-zero-divisor on \(\iota \to M\) iff it is one on
  \(M\) (\(\iota\) nonempty for the converse); and by \ref{pa-lem:ideal_smul_top_pi} at the
  ideal \((x)\) there is an \(R\)-linear equivalence
  \((\iota \to M)/x(\iota \to M) \cong (\iota \to M/xM)\) (a quotient of a finite
  product by a componentwise submodule is the product of the quotients), so by
  transport of regular sequences along linear equivalences (as in
  \ref{pa-lem:depth_linearEquiv}) and the inductive hypothesis applied to \(M/xM\), the
  tails match as well.
\end{proof}

\begin{lemma}[Minimal surjections from a free module]
  \label{pa-lem:minimal_surjection}
  \lean{RingTheory.Module.exists_minimalSurjection_finite_localRing}
  \leanok
  \dcref{custom}

  Let \((R, \mathfrak m)\) be a local ring with residue field \(\kappa\) and \(M\) a
  finite \(R\)-module. There exist \(n = \dim_\kappa(\kappa \otimes_R M)\) and a
  surjective \(R\)-linear map \(f : R^n \to M\) with
  \(\ker f \subseteq \mathfrak m \cdot R^n\).
\end{lemma}
\begin{proof}
  \leanok
  Choose a \(\kappa\)-basis \(b_1, \dots, b_n\) of \(\kappa \otimes_R M = M/\mathfrak m M\)
  and lift each \(b_i\) to \(m_i \in M\) along the surjection
  \(M \to \kappa \otimes_R M\), \(m \mapsto 1 \otimes m\). Let \(f : R^n \to M\) send
  the \(i\)-th standard basis vector to \(m_i\), so
  \(f(x) = \sum_i x_i m_i\). The images \(1 \otimes m_i = b_i\) span
  \(\kappa \otimes_R M\), so by Nakayama's lemma the \(m_i\) generate \(M\): \(f\) is
  surjective. If \(f(x) = 0\), applying \(1 \otimes -\) gives
  \(\sum_i \bar{x_i}\, b_i = 0\) in \(\kappa \otimes_R M\); linear independence of the
  \(b_i\) forces every \(\bar{x_i} = 0\), i.e.\ every \(x_i \in \mathfrak m\), so
  \(x \in \mathfrak m \cdot R^n\).
\end{proof}

\begin{lemma}[From the dimension equation to the vanishing bound]
  \label{pa-lem:pd_lt_of_pd_eq}
  \lean{RingTheory.Module.hasProjectiveDimensionLT_succ_of_projectiveDimension_eq}
  \leanok
  \dcref{custom}

  \uses{pa-def:projectiveDimension}
  Let \(R\) be a ring, \(M\) an \(R\)-module and \(n \in \mathbb{N}\) with
  \(\mathrm{pd}_R(M) = n\). Then \(M\) has projective dimension \(< n + 1\), i.e.\
  \(\Ext^i_R(M, -) = 0\) for all \(i \geq n + 1\).
\end{lemma}
\begin{proof}
  \leanok
  Immediate from Definition~\ref{pa-def:projectiveDimension}: the projective dimension is
  the infimum of the bounds \(m\) with \(\Ext^i(M,-) = 0\) for \(i > m\), and it equals
  \(n < n+1\).
\end{proof}

\begin{lemma}[Syzygy descent]
  \label{pa-lem:syzygy_descent}
  \lean{RingTheory.Module.hasProjectiveDimensionLT_ker_of_surjection}
  \leanok
  \dcref{custom}

  \uses{pa-lem:ext_covariant_les, pa-lem:ext_contravariant_les}
  Let \(R\) be a commutative ring, \(f : R^n \to M\) a surjective \(R\)-linear map, and
  \(k \in \mathbb{N}\). If \(M\) has projective dimension \(< k + 2\), then
  \(K = \ker f\) has projective dimension \(< k + 1\).
\end{lemma}
\begin{proof}
  \leanok
  Consider the short exact sequence \(0 \to K \to R^n \to M \to 0\) and any test module
  \(Y\). For \(i \geq k + 1\), the contravariant long exact sequence
  (\ref{pa-lem:ext_contravariant_les}) contains the exact segment
  \[
    \Ext^i(R^n, Y) \longrightarrow \Ext^i(K, Y) \longrightarrow \Ext^{i+1}(M, Y).
  \]
  The first term vanishes because \(R^n\) is free, hence projective, and \(i \geq 1\);
  the last vanishes because \(i + 1 \geq k + 2\) and \(M\) has projective dimension
  \(< k+2\). Hence \(\Ext^i(K, Y) = 0\) for all \(i \geq k+1\) and all \(Y\).
\end{proof}

\begin{lemma}[Syzygy ascent]
  \label{pa-lem:syzygy_ascent}
  \lean{RingTheory.Module.hasProjectiveDimensionLT_succ_of_hasProjectiveDimensionLT_ker}
  \leanok
  \dcref{custom}

  \uses{pa-lem:ext_contravariant_les}
  Let \(R\) be a commutative ring, \(f : R^n \to M\) a surjective \(R\)-linear map, and
  \(k \in \mathbb{N}\). If \(K = \ker f\) has projective dimension \(< k + 1\), then
  \(M\) has projective dimension \(< k + 2\).
\end{lemma}
\begin{proof}
  \leanok
  For a test module \(Y\) and \(i \geq k + 2\), the contravariant long exact sequence
  (\ref{pa-lem:ext_contravariant_les}) for \(0 \to K \to R^n \to M \to 0\) contains
  \[
    \Ext^{i-1}(K, Y) \xrightarrow{\ \delta\ } \Ext^{i}(M, Y) \longrightarrow
      \Ext^{i}(R^n, Y).
  \]
  The first term vanishes because \(i - 1 \geq k + 1\), the last because \(R^n\) is
  projective and \(i \geq 1\). Hence \(\Ext^i(M, Y) = 0\).
\end{proof}

\begin{lemma}[The two depth inequalities of the kernel sequence]
  \label{pa-lem:depth_ses_kernel}
  \lean{RingTheory.Module.depth_ses_ineqs_of_surjection_finite_localRing}
  \leanok
  \dcref{custom}

  \uses{pa-thm:depth_ses, pa-lem:depth_pi_const}
  Let \((R, \mathfrak m)\) be a Noetherian local ring, \(M\) a nonzero finite
  \(R\)-module, \(n \geq 1\), and \(f : R^n \to M\) a surjective \(R\)-linear map with
  nonzero kernel \(K\). Then
  \[
    \min\{\mathrm{depth}(R),\, \mathrm{depth}(K) - 1\} \leq \mathrm{depth}(M)
    \quad\text{and}\quad
    \min\{\mathrm{depth}(R),\, \mathrm{depth}(M) + 1\} \leq \mathrm{depth}(K).
  \]
\end{lemma}
\begin{proof}
  \leanok
  \(K\) is finite (a submodule of a finite module over a Noetherian ring), so parts (2)
  and (3) of \ref{pa-thm:depth_ses} apply to \(0 \to K \to R^n \to M \to 0\); substitute
  \(\mathrm{depth}(R^n) = \mathrm{depth}(R)\) from \ref{pa-lem:depth_pi_const}
  (\(n \geq 1\)).
\end{proof}

\begin{lemma}[A nonzero class at the depth index]
  \label{pa-lem:ext_nonzero_at_depth}
  \lean{RingTheory.Module.exists_ne_zero_ext_of_depth_eq}
  \leanok
  \dcref{custom}

  \uses{pa-thm:depth_ext}
  Let \((R, \mathfrak m)\) be a Noetherian local ring with residue field \(\kappa\) and
  \(M\) a nonzero finite \(R\)-module with \(\mathrm{depth}(M) = D \in \mathbb{N}\).
  Then \(\Ext^D_R(\kappa, M) \neq 0\).
\end{lemma}
\begin{proof}
  \leanok
  By \ref{pa-thm:depth_ext}, \(\Ext^i(\kappa, M) = 0\) for all \(i < D\), while
  \(D + 1 \leq \mathrm{depth}(M)\) fails, so the vanishing of all \(\Ext^i(\kappa, M)\)
  with \(i < D + 1\) fails; the failure can only occur at \(i = D\).
\end{proof}

\begin{lemma}[Matrix collapse on \(\Ext\)]
  \label{pa-lem:ext_matrix_collapse}
  \lean{RingTheory.Module.ext_comp_mk₀_ofHom_eq_zero_of_entries_mem_annihilator}
  \leanok
  \dcref{custom}

  \uses{pa-lem:ext_annihilator_smul}
  Let \(R\) be a commutative ring, \(N\) an \(R\)-module, and
  \(A : R^m \to R^n\) an \(R\)-linear map all of whose matrix entries
  \(a_{ij} = (A e_j)_i\) lie in \(\mathrm{Ann}_R(N)\). Then for every \(p \geq 0\) the
  postcomposition map \(A_* : \Ext^p_R(N, R^m) \to \Ext^p_R(N, R^n)\) is zero.
\end{lemma}
\begin{proof}
  \leanok
  Write \(A = \sum_{i,j} a_{ij}\, E_{ij}\), where \(E_{ij} : R^m \to R^n\) is the
  elementary map sending \(e_j\) to \(e_i\) and the other basis vectors to \(0\)
  (evaluate both sides on each \(e_j\)). Postcomposition on \(\Ext^p(N, -)\) is additive
  and \(R\)-linear in the map, so for \(e \in \Ext^p(N, R^m)\),
  \(A_*(e) = \sum_{i,j} a_{ij} \cdot (E_{ij})_*(e)\); each summand is the action of
  \(a_{ij} \in \mathrm{Ann}_R(N)\) on a class of \(\Ext^p(N, R^n)\), hence zero by
  \ref{pa-lem:ext_annihilator_smul}.
\end{proof}

\begin{theorem}[Auslander--Buchsbaum, positive projective dimension]
  \label{pa-thm:auslander_buchsbaum_succ}
  \lean{RingTheory.auslander_buchsbaum_formula_succ_pd}
  \leanok
  \source{stacks-project}

  \uses{pa-lem:minimal_surjection, pa-lem:pd_lt_of_pd_eq, pa-lem:syzygy_descent,
    pa-lem:syzygy_ascent, pa-lem:depth_ses_kernel, pa-lem:ext_nonzero_at_depth,
    pa-lem:ext_matrix_collapse, pa-lem:ext_covariant_les, pa-thm:depth_ext, pa-thm:depth_ses,
    pa-lem:depth_pi_const, pa-lem:ideal_smul_top_pi, pa-lem:depth_linearEquiv}
  \dcref{custom}
  Let \((R, \mathfrak m)\) be a Noetherian local ring and \(M\) a nonzero finite
  \(R\)-module with \(\mathrm{pd}_R(M) = k + 1\) for some \(k \in \mathbb{N}\). Then
  \[
    (k + 1) + \mathrm{depth}(M) = \mathrm{depth}(R).
  \]
\end{theorem}
\begin{proof}
  \leanok
  Induct on \(k\), generalising \(M\).

  \emph{Base case \(\mathrm{pd}_R(M) = 1\).} Take a minimal surjection
  \(f : R^n \twoheadrightarrow M\) with \(K = \ker f \subseteq \mathfrak m \cdot R^n\)
  (\ref{pa-lem:minimal_surjection}); \(n \geq 1\) since \(M \neq 0\). By
  \ref{pa-lem:pd_lt_of_pd_eq} and \ref{pa-lem:syzygy_descent} (at \(k = 0\)), \(K\) has
  projective dimension \(< 1\), i.e.\ \(K\) is projective; it is finite (Noetherian) and
  \(R\) is local, so \(K\) is free. \(K \neq 0\): otherwise \(f\) is an isomorphism,
  \(M\) is free and \(\mathrm{pd}_R(M) = 0 \neq 1\). Fix an isomorphism
  \(K \cong R^{k'}\) with \(k' \geq 1\) and let \(A : R^{k'} \to R^n\) be the composite
  inclusion. Every matrix entry of \(A\) lies in \(\mathfrak m\): the columns of \(A\)
  lie in \(K \subseteq \mathfrak m \cdot R^n\), which is componentwise
  \(\mathfrak m\)-valued by \ref{pa-lem:ideal_smul_top_pi}; and
  \(\mathfrak m = \mathrm{Ann}_R(\kappa)\) for the residue field \(\kappa\).

  \(\geq\): part (2) of \ref{pa-thm:depth_ses} applied to
  \(0 \to R^{k'} \xrightarrow{A} R^n \to M \to 0\), with
  \(\mathrm{depth}(R^{k'}) = \mathrm{depth}(R^n) = \mathrm{depth}(R)\)
  (\ref{pa-lem:depth_pi_const}), gives
  \(\mathrm{depth}(M) \geq \min\{\mathrm{depth}(R), \mathrm{depth}(R) - 1\}
  = \mathrm{depth}(R) - 1\), i.e.\
  \(\mathrm{depth}(R) \leq 1 + \mathrm{depth}(M)\).

  \(\leq\): if \(\mathrm{depth}(R) = \infty\) there is nothing to prove beyond the
  \(\geq\) direction forcing \(\mathrm{depth}(M) = \infty\); so let
  \(\mathrm{depth}(R) = D < \infty\). Pick a nonzero class
  \(\alpha \in \Ext^D(\kappa, R^{k'})\) (\ref{pa-lem:ext_nonzero_at_depth}, using
  \(\mathrm{depth}(R^{k'}) = D\)). By \ref{pa-lem:ext_matrix_collapse},
  \(A_*\alpha = 0\). If \(D = 0\), postcomposition with the monomorphism \(A\) is
  injective on \(\Ext^0\) (\ref{pa-lem:ext_covariant_les}), contradicting
  \(\alpha \neq 0\); so \(D = D' + 1\). The covariant long exact sequence
  (\ref{pa-lem:ext_covariant_les}) contains
  \[
    \Ext^{D'}(\kappa, M) \xrightarrow{\ \delta\ } \Ext^{D}(\kappa, R^{k'})
      \xrightarrow{\ A_*\ } \Ext^{D}(\kappa, R^n),
  \]
  so \(\alpha = \delta(\beta)\) for some \(\beta \in \Ext^{D'}(\kappa, M)\), and
  \(\beta \neq 0\). Hence \(\Ext^{D'}(\kappa, M) \neq 0\), so
  \(\mathrm{depth}(M) \leq D'\) by \ref{pa-thm:depth_ext}, i.e.\
  \(1 + \mathrm{depth}(M) \leq D = \mathrm{depth}(R)\). Together:
  \(1 + \mathrm{depth}(M) = \mathrm{depth}(R)\).

  \emph{Inductive step \(\mathrm{pd}_R(M) = k + 2\).} Take a minimal surjection
  \(f : R^n \twoheadrightarrow M\) (\ref{pa-lem:minimal_surjection}), \(n \geq 1\), and set
  \(K = \ker f\). Then \(\mathrm{pd}_R(K) = k + 1\) exactly: \(\leq\) by
  \ref{pa-lem:pd_lt_of_pd_eq} and \ref{pa-lem:syzygy_descent}; and if \(K\) had projective
  dimension \(< k + 1\), then \ref{pa-lem:syzygy_ascent} would give
  \(\mathrm{pd}_R(M) < k + 2\), a contradiction. In particular \(K \neq 0\) (the zero
  module has projective dimension \(-\infty\)), and \(K\) is finite. The inductive
  hypothesis applied to \(K\) gives
  \((k + 1) + \mathrm{depth}(K) = \mathrm{depth}(R)\), and \ref{pa-lem:depth_ses_kernel}
  gives the two inequalities
  \(\min\{\mathrm{depth}(R), \mathrm{depth}(K) - 1\} \leq \mathrm{depth}(M)\) and
  \(\min\{\mathrm{depth}(R), \mathrm{depth}(M) + 1\} \leq \mathrm{depth}(K)\).

  It remains to combine these in \(\{0, 1, \dots, \infty\}\). If
  \(\mathrm{depth}(K) = \infty\), the hypothesis forces
  \(\mathrm{depth}(R) = \infty\), the first inequality forces
  \(\mathrm{depth}(M) = \infty\), and both sides of the claim are \(\infty\). If
  \(\mathrm{depth}(K) = K_0 < \infty\) and \(\mathrm{depth}(M) = \infty\), the second
  inequality reads \(\min\{(k+1) + K_0, \infty\} = (k+1) + K_0 \leq K_0\), impossible
  since \(k + 1 \geq 1\). So both are finite, say \(\mathrm{depth}(M) = M_0\), and
  \(\mathrm{depth}(R) = (k+1) + K_0 > K_0 - 1\); the first inequality then reads
  \(K_0 - 1 \leq M_0\). In the second inequality the minimum cannot be
  \(\mathrm{depth}(R)\), since \((k+1) + K_0 \leq K_0\) is impossible; so it reads
  \(M_0 + 1 \leq K_0\). Hence \(M_0 = K_0 - 1\) with \(K_0 \geq 1\), and
  \((k + 2) + M_0 = (k + 1) + K_0 = \mathrm{depth}(R)\).
\end{proof}

\begin{theorem}[The Auslander--Buchsbaum formula]
  \label{pa-thm:auslander_buchsbaum}
  \lean{RingTheory.auslander_buchsbaum_formula}
  \leanok
  \source{stacks-project}

  \uses{pa-thm:auslander_buchsbaum_succ, pa-lem:depth_pi_const, pa-lem:depth_linearEquiv,
    pa-def:projectiveDimension}
  \dcref{custom}
  Let \((R, \mathfrak m)\) be a Noetherian local ring and \(M\) a nonzero finite
  \(R\)-module of finite projective dimension \(\mathrm{pd}_R(M) = n \in \mathbb{N}\).
  Then
  \[
    \mathrm{pd}_R(M) + \mathrm{depth}(M) = \mathrm{depth}(R).
  \]
\end{theorem}
\begin{proof}
  \leanok
  If \(n = 0\), then \(M\) is projective (Definition~\ref{pa-def:projectiveDimension}),
  finite over the local Noetherian ring \(R\), hence free: \(M \cong R^{k}\) with
  \(k = \mathrm{rank}\, M \geq 1\) (as \(M \neq 0\)). By \ref{pa-lem:depth_linearEquiv} and
  \ref{pa-lem:depth_pi_const}, \(\mathrm{depth}(M) = \mathrm{depth}(R^k) =
  \mathrm{depth}(R)\), which is the formula with \(n = 0\). If \(n = k + 1\), this is
  \ref{pa-thm:auslander_buchsbaum_succ}.
\end{proof}

\section{Cohen--Macaulay and regular local rings}
\label{pa-sec:regular_cm}

This section defines the Cohen--Macaulay condition for a Noetherian local ring and
proves the two classical structure theorems for regular local rings: a regular local
ring is a domain, and a regular system of parameters is a regular sequence --- so
regular local rings are Cohen--Macaulay. Throughout, for a Noetherian local ring
\((R, \mathfrak m)\) we write \(\mu(\mathfrak m)\) for the minimal number of generators
of \(\mathfrak m\), which by Nakayama's lemma equals
\(\dim_\kappa \mathfrak m/\mathfrak m^2\), the dimension of the cotangent space; recall
that \(R\) is \emph{regular} when \(\mu(\mathfrak m) = \dim R\), the Krull dimension
(and that \(\dim R \leq \mu(\mathfrak m)\) always, by Krull's height theorem).

\begin{definition}[Cohen--Macaulay local ring]
  \label{pa-def:cohen_macaulay}
  \lean{RingTheory.CohenMacaulay}
  \leanok
  \source{stacks-project}

  \uses{pa-def:depth}
  \dcref{custom}
  A Noetherian local ring \((R, \mathfrak m)\) is \emph{Cohen--Macaulay} if
  \(\mathrm{depth}(R) = \dim R\), its depth (Definition~\ref{pa-def:depth}) equals its
  Krull dimension.
\end{definition}

\begin{theorem}[Cotangent dimension drop]
  \label{pa-thm:cotangent_dim_drop}
  \lean{RingTheory.CohenMacaulay.finrank_cotangentSpace_quot_span_singleton_succ}
  \leanok
  \dcref{custom}

  Let \((R, \mathfrak m)\) be a Noetherian local ring and
  \(x \in \mathfrak m \setminus \mathfrak m^2\). Write \(\mathfrak m'\) for the maximal
  ideal of the local ring \(R/(x)\). Then
  \[
    \dim_{\kappa'} \mathfrak m'/\mathfrak m'^2 + 1
      \;=\; \dim_{\kappa} \mathfrak m/\mathfrak m^2,
  \]
  where \(\kappa, \kappa'\) are the residue fields; equivalently
  \(\mu(\mathfrak m') + 1 = \mu(\mathfrak m)\).
\end{theorem}
\begin{proof}
  \leanok
  By Nakayama both sides may be computed as minimal generator counts, so we prove
  \(\mu(\mathfrak m') + 1 = \mu(\mathfrak m)\), by two inequalities.

  \(\leq\) (Steinitz exchange): let \(V = \{v_1, \dots, v_d\}\) be a generating set of
  \(\mathfrak m\) with \(d = \mu(\mathfrak m)\); note \(d \geq 1\), since \(d = 0\)
  would give \(\mathfrak m = 0\) and \(x \in \mathfrak m^2\). Write
  \(x = \sum_i c_i v_i\). If every \(c_i \in \mathfrak m\), then
  \(x \in \mathfrak m \cdot \mathfrak m = \mathfrak m^2\), a contradiction; so some
  \(c_{i_0}\) is a unit, and
  \(v_{i_0} = c_{i_0}^{-1}\bigl(x - \sum_{i \neq i_0} c_i v_i\bigr)\) lies in the span
  of \((V \setminus \{v_{i_0}\}) \cup \{x\}\). Hence
  \((V \setminus \{v_{i_0}\}) \cup \{x\}\) still generates \(\mathfrak m\); its image in
  \(R/(x)\) generates \(\mathfrak m' = \mathfrak m/(x)\), and the image of \(x\) is
  zero, so the \(d - 1\) images of \(V \setminus \{v_{i_0}\}\) generate
  \(\mathfrak m'\): \(\mu(\mathfrak m') \leq d - 1\).

  \(\geq\) (lift and cons): let \(T\) be a generating set of \(\mathfrak m'\) of size
  \(\mu(\mathfrak m')\) and lift each element along the surjection
  \(R \twoheadrightarrow R/(x)\) to \(\widetilde T \subseteq \mathfrak m\) (the
  preimage of \(\mathfrak m'\) is \(\mathfrak m\), the unique maximal ideal). Since
  \(\mathfrak m\) is the preimage of \(\mathfrak m' = (\text{image of } \widetilde T)\)
  and the kernel of the surjection is \((x)\), we get
  \(\mathfrak m = (\widetilde T) + (x)\), so \(\widetilde T \cup \{x\}\) generates
  \(\mathfrak m\): \(\mu(\mathfrak m) \leq \mu(\mathfrak m') + 1\).
\end{proof}

\begin{lemma}[The Nakayama witness \(x \in \mathfrak m \setminus \mathfrak m^2\)]
  \label{pa-lem:notMemSq_witness}
  \lean{RingTheory.CohenMacaulay.exists_notMemSq_of_spanFinrank_pos}
  \leanok
  \dcref{custom}

  Let \((R, \mathfrak m)\) be a Noetherian local ring with \(\mu(\mathfrak m) \geq 1\).
  Then there exists \(x \in \mathfrak m\) with \(x \notin \mathfrak m^2\).
\end{lemma}
\begin{proof}
  \leanok
  Otherwise \(\mathfrak m \subseteq \mathfrak m^2 = \mathfrak m \cdot \mathfrak m\);
  since \(\mathfrak m\) is finitely generated (Noetherian) and contained in the
  Jacobson radical, Nakayama's lemma gives \(\mathfrak m = 0\), so
  \(\mu(\mathfrak m) = 0\) --- a contradiction.
\end{proof}

\begin{lemma}[The regular quotient step]
  \label{pa-lem:regular_quotient}
  \lean{RingTheory.CohenMacaulay.regularLocal_quotient_isRegularLocal_of_notMemSq}
  \leanok
  \dcref{custom}

  \uses{pa-thm:cotangent_dim_drop, pa-lem:notMemSq_witness}
  Let \(R\) be a regular local Noetherian ring with \(\mu(\mathfrak m) = k + 1\) and let
  \(x \in \mathfrak m \setminus \mathfrak m^2\). Then \(R/(x)\) is a nonzero regular
  local Noetherian ring with \(\mu(\mathfrak m') = k\), where \(\mathfrak m'\) is its
  maximal ideal. (No non-zero-divisor hypothesis on \(x\) is made.)
\end{lemma}
\begin{proof}
  \leanok
  Since \(x \in \mathfrak m\) is a nonunit, \((x) \neq R\) and \(R/(x)\) is a nonzero
  local Noetherian ring. By \ref{pa-thm:cotangent_dim_drop},
  \(\mu(\mathfrak m') = k\). For the Krull dimension: since \(x\) lies in the Jacobson
  radical of \(R\), quotienting by the single element \(x\) can drop the dimension by at
  most one, \(\dim R \leq \dim R/(x) + 1\) (the general bound
  \(\dim R \leq \dim R/(S) + \#S\) for a subset \(S\) of the Jacobson radical, which
  does \emph{not} require \(x\) to be a non-zero-divisor). Regularity of \(R\) gives
  \(\dim R = \mu(\mathfrak m) = k + 1\), so \(\dim R/(x) \geq k\); Krull's height
  theorem gives \(\dim R/(x) \leq \mu(\mathfrak m') = k\). Hence
  \(\dim R/(x) = k = \mu(\mathfrak m')\) and \(R/(x)\) is regular.
\end{proof}

\begin{theorem}[Regular local rings are domains]
  \label{pa-thm:regular_isDomain}
  \lean{RingTheory.CohenMacaulay.isDomain_of_regularLocal}
  \leanok
  \source{stacks-project}

  \uses{pa-lem:notMemSq_witness, pa-lem:regular_quotient}
  \dcref{custom}
  Every regular local Noetherian ring is an integral domain.
\end{theorem}
\begin{proof}
  \leanok
  Strong induction on \(d = \mu(\mathfrak m)\), generalising the ring.

  \emph{Base case \(d = 0\).} Then \(\mathfrak m = 0\), so \(R\) is a field, hence a
  domain.

  \emph{Inductive step \(d = k + 1\).} Pick \(x \in \mathfrak m \setminus \mathfrak m^2\)
  (\ref{pa-lem:notMemSq_witness}). By \ref{pa-lem:regular_quotient}, \(R/(x)\) is regular
  local with \(\mu = k\), hence a domain by the inductive hypothesis; so \((x)\) is a
  prime ideal. Choose a minimal prime \(\mathfrak p \subseteq (x)\) of \(R\) (every
  prime contains a minimal prime).

  \emph{Case \(x \notin \mathfrak p\).} For \(y \in \mathfrak p \subseteq (x)\) write
  \(y = xz\); since \(\mathfrak p\) is prime and \(x \notin \mathfrak p\), we get
  \(z \in \mathfrak p\). Hence \(\mathfrak p = x\,\mathfrak p\) with
  \(x\) in the Jacobson radical and \(\mathfrak p\) finitely generated; Nakayama's lemma
  gives \(\mathfrak p = 0\). A minimal prime equal to \((0)\) means \((0)\) is prime, so
  \(R\) is a domain.

  \emph{Case \(x \in \mathfrak p\).} Then \((x) \subseteq \mathfrak p \subseteq (x)\),
  so \((x)\) itself is a minimal prime of \(R\). We refute this. \(R\) is Noetherian, so
  it has only finitely many minimal primes. The maximal ideal \(\mathfrak m\) is
  contained in none of the ideals \(\mathfrak m^2, \mathfrak q_1, \dots, \mathfrak q_r\)
  (\(\mathfrak q_i\) the minimal primes): containment in \(\mathfrak m^2\) would give
  \(\mathfrak m = 0\) by Nakayama, contradicting \(d \geq 1\); and
  \(\mathfrak m \subseteq \mathfrak q_i\) would force \(\mathfrak q_i = \mathfrak m\)
  (primes of a local ring lie in \(\mathfrak m\)), making \(\mathfrak m\) a minimal
  prime, so \(\dim R = 0\) --- but \(\dim R = \mu(\mathfrak m) = k + 1 \geq 1\) by
  regularity. By prime avoidance (with the single non-prime member \(\mathfrak m^2\)
  allowed), \(\mathfrak m \not\subseteq \mathfrak m^2 \cup \mathfrak q_1 \cup \dots \cup
  \mathfrak q_r\); pick \(x'\) in \(\mathfrak m\) outside this union. Then
  \(x' \in \mathfrak m \setminus \mathfrak m^2\) avoids every minimal prime. Run the
  preceding two paragraphs with \(x'\) in place of \(x\): \(R/(x')\) is a domain by
  \ref{pa-lem:regular_quotient} and the inductive hypothesis, so \((x')\) is prime; a
  minimal prime \(\mathfrak p' \subseteq (x')\) cannot contain \(x'\), so the previous
  case applies and shows \(\mathfrak p' = 0\) and \(R\) is a domain. But in a domain the
  only minimal prime is \((0)\), so the minimal prime \((x)\) equals \((0)\), giving
  \(x = 0 \in \mathfrak m^2\) --- contradicting \(x \notin \mathfrak m^2\).
\end{proof}

\begin{theorem}[A regular system of parameters is a regular sequence]
  \label{pa-thm:regular_system_isRegular}
  \lean{RingTheory.CohenMacaulay.exists_isRegular_of_regularLocal}
  \leanok
  \dcref{custom}

  \uses{pa-lem:notMemSq_witness, pa-thm:regular_isDomain, pa-lem:regular_quotient}
  Let \(R\) be a regular local Noetherian ring. Then there is an \(R\)-regular sequence
  contained in \(\mathfrak m\) of length \(\mu(\mathfrak m) = \dim R\).
\end{theorem}
\begin{proof}
  \leanok
  Induct on \(d = \mu(\mathfrak m)\), generalising the ring. For \(d = 0\) take the
  empty sequence (\(R \neq 0\), so it is vacuously \(R\)-regular). For \(d = k + 1\):
  pick \(x \in \mathfrak m \setminus \mathfrak m^2\) (\ref{pa-lem:notMemSq_witness}); then
  \(x \neq 0\) (zero lies in \(\mathfrak m^2\)) and \(R\) is a domain
  (\ref{pa-thm:regular_isDomain}), so \(x\) is a non-zero-divisor on \(R\). By
  \ref{pa-lem:regular_quotient}, \(R/(x)\) is regular local with
  \(\mu(\mathfrak m') = k\), so by the inductive hypothesis there is an
  \(R/(x)\)-regular sequence \(\bar y_1, \dots, \bar y_k\) in \(\mathfrak m'\). Lift
  each \(\bar y_i\) to \(y_i \in \mathfrak m\) (the preimage of \(\mathfrak m'\) is
  \(\mathfrak m\)). Regularity of a sequence on the \(R\)-module \(R/xR\) is computed
  through the images of its entries in \(R/(x)\), so \(y_1, \dots, y_k\) is a regular
  sequence on the \(R\)-module \(R/xR\); consing, \(x, y_1, \dots, y_k\) is an
  \(R\)-regular sequence of length \(k + 1\) in \(\mathfrak m\).
\end{proof}

\begin{theorem}[Regular local rings are Cohen--Macaulay]
  \label{pa-thm:regular_is_CM}
  \lean{RingTheory.CohenMacaulay.of_regular}
  \leanok
  \source{stacks-project}

  \uses{pa-def:cohen_macaulay, pa-def:depth, pa-thm:regular_system_isRegular}
  \dcref{custom}
  Every regular local Noetherian ring \(R\) is Cohen--Macaulay:
  \(\mathrm{depth}(R) = \dim R\).
\end{theorem}
\begin{proof}
  \leanok
  The side condition of Definition~\ref{pa-def:depth} is in the supremum branch:
  \(\mathfrak m \cdot R = \mathfrak m \neq R\).

  \emph{Lower bound.} By \ref{pa-thm:regular_system_isRegular} there is an \(R\)-regular
  sequence in \(\mathfrak m\) of length \(\mu(\mathfrak m) = \dim R\), so
  \(\mathrm{depth}(R) \geq \dim R\).

  \emph{Upper bound.} Every \(R\)-regular sequence \(x_1, \dots, x_\ell\) (in any
  ideal) satisfies \(\ell \leq \dim R\): for a Noetherian local ring, quotienting by a
  regular sequence drops the Krull dimension by exactly its length,
  \(\dim R/(x_1, \dots, x_\ell) + \ell = \dim R\), and the quotient is nonzero (the
  nondegeneracy condition \((x_1, \dots, x_\ell)R \neq R\)), so its dimension is
  \(\geq 0\). Hence the supremum defining \(\mathrm{depth}(R)\) is at most
  \(\dim R\).

  Therefore \(\mathrm{depth}(R) = \dim R = \mu(\mathfrak m)\), and \(R\) is
  Cohen--Macaulay (Definition~\ref{pa-def:cohen_macaulay}).
\end{proof}
